%%%%%%%%%%%%%%%%%%%%%%%%%%%%%%%%%%%%%%%%%%%%%%%%%%%%%%%%%%%%%%%%%%%%%%%%%%%%%
%%%%%%                                                                  %%%%% 
%%%%%%          Maqueta de memòria TFC/PFC de l'EETAC                   %%%%% 
%%%%%%                                                                  %%%%% 
%%%%%%%%%%%%%%%%%%%%%%%%%%%%%%%%%%%%%%%%%%%%%%%%%%%%%%%%%%%%%%%%%%%%%%%%%%%%%
%%%%%%%%%%%%%%%%%%%%%%%%%%%%%%%%%%%%%%%%%%%%%%%%%%%%%%%%%%%%%%%%%%%%%%%%%%%%%
%%                                                                         %%
%%          Autor: Xavier Prats i Menéndez (xavier.prats@upc.edu)          %% 
%%                  Technical University of Catalonia (UPC)                %%
%%                                                                         %%
%%%%%%%%%%%%%%%%%%%%%%%%%%%%%%%%%%%%%%%%%%%%%%%%%%%%%%%%%%%%%%%%%%%%%%%%%%%%%
%%      This work is licensed under the Creative Commons  Attribution-     %%
%%   -Noncommercial-Share Alike 3.0 Spain License. To view a copy of this  %% 
%%    license, visit http://creativecommons.org/licenses/by-nc-sa/3.0/es/  %%
%%    or send a letter to Creative Commons, 171 Second Street, Suite 300,  %%
%%                  San Francisco,California, 94105, USA.                  %%
%%%%%%%%%%%%%%%%%%%%%%%%%%%%%%%%%%%%%%%%%%%%%%%%%%%%%%%%%%%%%%%%%%%%%%%%%%%%%
%% Versió 2.1 - Juliol 2012                                                %%
%%%%%%%%%%%%%%%%%%%%%%%%%%%%%%%%%%%%%%%%%%%%%%%%%%%%%%%%%%%%%%%%%%%%%%%%%%%%%

%%% NOTA: els seguents packages son necessaris per utilitzar la
%%%       plantilla seguent:
%%%       ifthen,calc,helvet,pslatex,fancyhdr,nextpage,subfigure,tocloft,graphicx,url

%%% NOTA: Es possible que algunes distribuicions Linux o Windows.
%%%       no portin aquests paquets instal·lats per defecte.
%%%       En aquest cas els haureu d'instal·lar manualment.


%%%%%%%%%%%%%%%%%%%%%%%%%%%%%%%%%%%%%%%%%%%%%%%%%%%%%%%%%%%%%%%%%%%%%%%%%%%%%
% 1- INICIALITZACIÓ
%%%%%%%%%%%%%%%%%%%%%%%%%%%%%%%%%%%%%%%%%%%%%%%%%%%%%%%%%%%%%%%%%%%%%%%%%%%%%

\documentclass[spanish,final]{setup/eetac_tfc_pfc}
%% * OPCIONS A CONFIGURAR al \documentclass
%%    - Estat del document: final o draft
%%      NOTA: Draft no inserta les figures i marca només l'espai que
%%      ocupen. També s'indica quan el text sobrepassa els marges.
%%      Draft és molt útil per compilar ràpid el document si no és important
%%      en aquell moment visualitzar les figures.
%%    - Idioma PRINCIPAL del document: catalan, spanish, english, french...

\usepackage[spanish]{babel}
%%  * INCLOURE TOTS ELS IDIOMES QUE S'USARAN EN EL DOCUMENT
%%    NOTA: per canviar d'idioma al mig del document usar:
%%          \selectlanguage{nom_idioma}
%%%%%%%%%%%%%%%%%%%%%%%%%%%%%%%%%%%%%%%%%%%%%%%%%%%%%%%%%%%%%%%%%%%%%%%%%%%%%

%%%%%%%%%%%%%%%%%%%%%%%%%%%%%%%%%%%%%%%%%%%%%%%%%%%%%%%%%%%%%%%%%%%%%%%%%%%%%
% 2- CÀRREGA DE PAQUETS ADICIONALS (OPCIONALS)
%%%%%%%%%%%%%%%%%%%%%%%%%%%%%%%%%%%%%%%%%%%%%%%%%%%%%%%%%%%%%%%%%%%%%%%%%%%%%

%%% NOTA: Es possible que algunes distribuicions Linux o Windows.
%%%       no portin aquests paquets instal·lats per defecte.
%%%       En aquest cas els haureu d'instal·lar manualment.

%% El paquet inputenc és extramadament útil. 
%% Permet escriure els accents directament amb l'editor de texte
%% sense haver de fer coses com per exemple: introducci\'o
%% Heu d'especificar la codificació de caracters que utilitzeu pel
%% vostre fitxer (en aquest exemple utf8)
\usepackage[normalem]{ulem}
\usepackage{xcolor}
\usepackage{matlab-prettifier}
\usepackage{cancel}
\definecolor{grey}{rgb}{0.5, 0.5, 0.52} 
\definecolor{navy}{rgb}{0.1, 0.05, 0.72} 
\newcommand{\PGPremove}[2]{\textcolor{red}{\sout{#1}}{ \textcolor{red}{#2}}}
\newcommand{\PGP}[1]{\textcolor{blue}{#1}}
\newcommand{\jordi}[1]{\textcolor{magenta}{#1}}
\newcommand{\sara}[1]{\textcolor{grey}{#1}}
\newcommand{\PGPcomment}[1]{\textcolor{navy}{#1}}
\usepackage[utf8]{inputenc}
\usepackage{float}

%% Símbols matemàtics de la American Mathematical Society
\usepackage{amssymb,amsmath,amsfonts,mathtools}  

%% El paquet array proporciona eines molt útils a l'hora de fer 
%% equacions amb matrius
\usepackage{array}             

%% Paquet que permet fer taules fusionant cel·les de files consecutives
\usepackage{multirow}  

%% Paquet molt útil en cas de tenir taules molt llargues que 
%   ocupin vàries pàgines
\usepackage{longtable}          

%% Permet canviar els colors del document
%\usepackage{color,colortbl}

%% Paquet molt útil que permet activar links en el PDF final.
%% * NO OBLIDAR DE CONFIGURAR els quatre primer camps!
\usepackage[
  pdfauthor={Dídac Bosch Garcias},            % Configurar adientment
  pdftitle={Trabajo fin de Grado - Dídac Bosch Garcias}, % Configurar adientment
  pdfsubject={Sistema de control de drones vía Flutter web},          % Configurar adientment  
  pdfkeywords={dron, Flutter, Ardupilot ...},    % Configurar adientment
  pdfcreator={EETAC-UPC}, 
  pdfproducer={LaTeX, dvipdf},
  pdfdisplaydoctitle=true, plainpages=false, linktocpage=true,         
  colorlinks=true, linkcolor=blue,citecolor=blue,urlcolor=blue,
  hyperfootnotes=false, pagebackref=true, pdfpagelabels=true,
  pdfpagemode=UseOutlines,
]{hyperref} 

%% NOTA IMPORTANT!:
%% Per tal que hyperef funcioni correctament amb els capitols o seccions no
%% numerats (\chapter*{}), com per exemple introducció, conclusions i bibliografia
%% cal posar les dues comandes seguents ABANS del \chapter*{} en questió
%\cleardoublepage
%\phantomsection

%% Permet trencar links URL. 
%% Atenció! afegir aquest paquet DESPRES del hyperref!!
\usepackage{breakurl} 

%% Permet arranjar matricialment multiples figures
%% NOTA: afegir aquest paquet DESPRES del hyperref!!
%%       Si no es desitja utilitzar aquest paquet, comentar la linia seguent
%%       i anar TAMBE al fitxer de classe (eetac_tfc_pfc.cls) per substituir: 
%%       \RequirePackage[subfigure]{tocloft}  per  \RequirePackage{tocloft}
\usepackage{subfigmat}       
\usepackage{acronym}

% ── Corrección de sangrías del ÍNDICE GENERAL ──────────────────
% Estilo limpio: número sin "CAPÍTULO", igual que la maqueta EETAC
\renewcommand{\cftchappresnum}{}          % quita "CAPÍTULO" del TOC
\renewcommand{\cftchapaftersnum}{.~~}     % añade punto y espacio
\setlength{\cftchapnumwidth}{2em}         % "1. " cabe en 2em

% Secciones (1.1, 1.2...) — sangría bajo el número de capítulo
\setlength{\cftsecindent}{2em}
\setlength{\cftsecnumwidth}{2.5em}

% Subsecciones (1.1.1, 1.2.3...) — sangría acumulada
\setlength{\cftsubsecindent}{4.5em}
\setlength{\cftsubsecnumwidth}{3.5em}

% Mostrar también subsubsecciones si las tienes
\setcounter{tocdepth}{3}

%%%%%%%%%%%%%%%%%%%%%%%%%%%%%%%%%%%%%%%%%%%%%%%%%%%%%%%%%%%%%%%%%%%%%%%%%%%%%


%%%%%%%%%%%%%%%%%%%%%%%%%%%%%%%%%%%%%%%%%%%%%%%%%%%%%%%%%%%%%%%%%%%%%%%%%%%%%
% 3- DOCUMENT
%%%%%%%%%%%%%%%%%%%%%%%%%%%%%%%%%%%%%%%%%%%%%%%%%%%%%%%%%%%%%%%%%%%%%%%%%%%%%

%%% Configuració de les dades i variables boleanes rellevants del document:
\input{setup/dades.tex}  

%%% Configuració de MACROS o ENTORNS (opcionals) definides per l'usuari:
\input{setup/user-macros.tex}  

%%% Configuració manual de les regles d'hyphenation:
\input{setup/hyphenation.tex}  

\begin{document}

%% Seleccionar l'idioma principal del document:
\selectlanguage{spanish}

\beforepreface  

%% RESUM i OVERVIEW
%%%%%%%%%%%%%%%%%%%%%%%%%%%%%%%%%%%%%%%%%%%%%%%%%%%%%%%%%%%%%%%%%%%%%%%%%%%%%
% NOTA: les longituds passades com a parametres d'entrada  s'han
%        d'ajustar manualment fins que el requadre del resum/overview
%        ocupi tota la pàgina. 

%%% Resum en català (o castellà)
\selectlanguage{spanish}   
\begin{resum}{1cm}
Abstract en català.

\end{resum}

%%% Resum en anglès
\selectlanguage{english}   
\begin{overview}{1cm}
Abstract in English

\end{overview}

% Tornar a l'idioma principal del document
\selectlanguage{spanish}  

%NOTA: En cas d'utilitzar l'espanyol com a idioma principal del document, el
%      latex anomena les taules com a 'Cuadros'. Si es desitja canviar aquesta
%      nomenclatura i utilitzar la paraula 'Tabla' descomentar les línies següents:
\def\listtablename{Índice de tablas}
\def\tablename{Tabla}%



% Amb aqueta comanda indiquem que ja s'han inclòs tots els apartats del prefaci del 
% projecte o podem començar a incloure els capitols de la memòria
\afterpreface


%%%%%%%%%%%%%%%%%%%%%%%%%%%%%%%%%%%%%%%%%%%%%%%%%%%%%%%%%%%%%%%%%%%%%%%%%%
%%%%%% INCLOURE A PARTIR D'AQUÍ TOTS ELS CAPÍTOLS DE LA MEMORIA   %%%%%%%%
%%%%%%%%%%%%%%%%%%%%%%%%%%%%%%%%%%%%%%%%%%%%%%%%%%%%%%%%%%%%%%%%%%%%%%%%%%

% NOTA: recordar que la introducció i les conclusions són capítols NO
%       enumerats, per tant s'ha d'usar \chapter*

% NOTA: és aconsellable incloure els capítols de la memòria en fitxers 
%       separats utlitzant la comanda \input  Per exemple:
%       \input{capitol1}  
%       que farà que s'inclogui el fitxer capitol1.tex

% NOTA: Si es vol incloure agraïments i/o glosari, fer-ho utilitzant 
% \chapter*{} i incloure'ls abans la introducció

\input{memoria/chapters/acronyms}
% chap1_introduccion.tex

\chapter{Introducción}   % capítulos numerados normales

\section{Contexto y motivación}
Los vehículos aéreos no tripulados (UAV, del inglés Unmanned Aerial Vehicle) han experimentado un crecimiento exponencial durante las últimas décadas, evolucionando desde aplicaciones militares, hasta convertirse en una herramienta de uso generalizado en sectores  civiles. comerciales y de investigación. En la actualidad, sus aplicaciones abarcan la agricultura de precisión, la inspección de infraestructuras, las operaciones de búsqueda y rescate y la cinematografía entre muchas otras\cite{UAV applications}.

El elemento central de cualquier operación con UAV es la estación de control en tierra (GCS, Ground Control Station), la interfaz software a través de la cual el operador monitoriza el estado del vehículo y controla el vuelo.  Históricamente, las GCS se han desarrollado como aplicaciones de escritorio fuertemente vinculadas a un sistema operativo concreto. Herramientas como Mission Planner\cite{missionplanner} y QGroundControl\cite{qgroundcontrol} representan el estándar de facto en código abierto, pero presentan una barrera de accesibilidad significativa. Requieren instalación local, no están preparadas para abrirse en web o móvil, y presentan una curva de aprendizaje pronunciada para operadores con poca experiencia.

A medida que las tecnologías web y móvil han crecido, ha surgido la oportunidad de construir interfaces de control ligeras y multiplataforma, que se ejecuten directamente en el navegador sin necesidad de instalación local. En contextos operativos, la posibilidad de conectarse a un dron desde cualquier dispositivo con navegador, o de presentar una interfaz sencilla de usar y entender, supone una mejora sustancial. 

Flutter\cite{flutter}, el framework de interfaz de usuario de código abierto desarrollado por Google, es un candidato adecuado para este tipo de aplicaciones. Permite apuntar a Android, iOS, web y escritorio, desde una única base de código, y su modelo de gestión de estados lo hace idóneo para interfaces que deben responder de forma continua a flujos de telemetría y reflejar cambios de estado del vehículo en tiempo real. A pesar de su creciente adopción en proyectos de IoT, existe una gran escasez de ejemplos documentados y reutilizables orientados específicamente al control de UAV. Los desarrolladores que deseen construir este tipo de interfaces tienen que crear soluciones a partir de ejemplos dispersos de dominios no relacionados, dificultando así los retos propios del control de vehículos en tiempo real.

Este proyecto nace para cubrir precisamente ese vacío mediante dos contribuciones complementarias, desarrolladas en el marco del DEE (Drone Engineering Ecosystem)\cite{dronsEETAC}. La primera es una colección de aplicaciones Flutter de demostración, independientes entre sí, cada una explorando un paradigma de control de dron específico: control RC mediante un doble joystick virtual, reconocimiento de comandos por voz, control inercial mediante la IMU del propio dispositivo, localización GPS del operador, streaming de vídeo con detección de objetos, planificación de misiones autónomas y registro y análisis de vuelos. Cada demostración está documentada de forma autónoma y es accesible públicamente como referencia para futuros desarrolladores del ecosistema.

La segunda contribución es EZDrone, una estación de control en tierra basada en Flutter Web que integra todos estos paradigmas en una arquitectura coherente. EZDrone permite controlar dos plataformas de vuelo presentes en el DEE. El dron principal, un chasis HexSon con controlador de vuelo Pixhawk\cite{pixhawk} ejecutando firmware ArduPilot\cite{ardupilot}, utilizado para operaciones en exteriores, y el dron DJI Tello\cite{tello}, usado para interiores, especialmente en demostraciones a visitantes. 

En conjunto, este proyecto amplía el bloque de aplicaciones Flutter del DEE con recursos reutilizables y una estación tierra web accesible desde cualquier dispositivo y cualquier navegador, orientada tanto al ecosistema abierto ArduPilot como al uso docente y demostrativo de la EETAC.



\section{Estado del arte}
\subsection{Estaciones de control en tierra de código abierto}
El ecosistema de UAV de código abierto está estructurado en torno a dos familias de autopilot, ArduPilot y PX4\cite{px4},  que emplean MAVLink\cite{mavlink} como protocolo de comunicación principal. Sobre este sustrato se han construido las principales GCS disponibles hoy en día.

Mission Planner es la  GCS de referencia para ArduPilot, desarrollada en C\#. Cubre el ciclo operativo completo: configuración de parámetros, calibración de sensores, planificación de misiones, monitorización de telemetría en tiempo real, control completo del UAV y análisis de logs post-vuelo. Es la herramienta elegida por la mayoría de los usuarios de ArduPilot, sin embargo, está limitada exclusivamente a Windows, y no contempla ninguna forma de despliegue web o móvil.

QGroundControl amplía significativamente el soporte de plataformas, es compatible con Windows, macOS, Linux, Android e iOS. Ofrece una interfaz más moderna y accesible. A pesar de sus versiones móviles, su funcionalidad está muy reducida respecto a la versión de escritorio, y no existe ninguna versión web.

La documentación oficial de ArduPilot identifica al menos diez GCS, pero ninguna de las opciones ofrece actualmente despliegue web\cite{ardupilot_gcs}. Esta ausencia es la brecha directa que EZDrone pretende cubrir.   


\subsection{Aplicaciones móviles para el control de drones}
En el sector comercial, fabricantes como DJI y Skydio han desarrollado aplicaciones propias que ofrecen una experiencia de usuario muy cuidada e integración completa con su hardware. Aplicaciones como DJI Fly\cite{djifly} presentan una interfaz simple y sencilla que cuenta con funciones de control de vuelo, transmisión de vídeo en tiempo real, modos de vuelo automatizados (como seguimiento de personas, vuelos orbitales, panorámicas) y edición de contenido directamente en el dispositivo. Skydio\cite{skydio}, por su parte, pone el foco en la autonomía de vuelo y la evasión de obstáculos, ofreciendo interfaces pensadas para que el operador pueda realizar vuelos complejos sin requerir de mucha experiencia y con la mínima intervención manual. Estas soluciones están muy pulidas desde el punto de vista de la experiencia de uso, es decir, tienen interfaces muy intuitivas y guiadas, y ofrecen conectividad fiable y funcional con los propios dispositivos del fabricante.

Sin embargo, este nivel de integración tiene un coste importante, es código cerrado, diseñado exclusivamente para su propio hardware. No existe ninguna interfaz de programación pública para adaptarlas a otros autopilots, protocolos de comunicación o plataformas de terceros. Esto las hace irrelevantes en entornos académicos como es el caso del DEE, donde la flexibilidad y reutilización del código son requisitos fundamentales.

Este proyecto, con EZDrone y las demostraciones, adopta un enfoque totalmente distinto. Al estar construido sobre Flutter Web y formar parte del DEE, es independiente del hardware (quitando el caso del Tello en EZDrone). Puede conectarse a cualquier plataforma que ejecute firmware ArduPilot, o extenderse para soportar otros autopilots. Las demostraciones son proyectos Flutter autónomos y documentados, pensados para ser estudiados, modificados e integrados para futuros desarrolladores. la cotrapartida es que el nivel de pulido visual es mucho menor que en una solución comercial dedicada. EZDrone está orientada más al control funcional, la docencia e investigación, no tanto a la experiencia de usuario del gran público. Esta diferencia de objetivos es la que justifica y delimita el enfoque adoptado en este trabajo.


\subsection{Flutter en IoT y robótica}
Flutter ha consolidado su presencia en el desarrollo IoT desde su lanzamiento en 2018\cite{flutter}, principalmente en cuadros de mando de sensores y aplicaciones de domótica que se comunican mediante MQTT. La combinación Flutter + MQTT\cite{mqtt_dart} cuenta con documentación razonable para casos de uso sencillos: lectura de valores de sensores, visualización de estado y envío de comandos discretos.

Su uso para el control de vehículos en tiempo real introduce retos cualitativamente distintos que están mucho menos documentados. El control de un UAV exige generar salidas de forma continua (como los valores de joystick enviados), gestionar múltiples flujos asíncronos simultáneamente (telemetría, vídeo, comandos de estado) y mantener la coherencia de la interfaz ante cambios de estado rápidos del vehículo. Trabajos recientes demuestran que arquitecturas similares (Flutter como frontend, MQTT para señalización y WebRTC para flujos de alta frecuencia) son técnicamente viables\cite{sun2025}. 

\section{TFG y TFM relacionados}
El control de drones mediante interfaces móviles y web, así como el desarrollo de aplicaciones Flutter, han sido objeto de Trabajos de Fin de Grado y Trabajos de Fin de Máster en los últimos años. Estos trabajos han explorado aspectos como la implementación de frontends para el control de UAV 

Este proyecto se enmarca en esa línea de trabajos, pero con un foco específico y diferenciado, la creación de una colección de demostraciones Flutter reutilizables que cubran los principales paradigmas de control de dron, y el desarrollo de EZDrone como estación de control web integrada en el DEE. 


\section{Estructura del documento}
El Capítulo 2 presenta el Drone Engineering Ecosystem (DEE), la plataforma académica desarrollada en la EETAC dentro de la cual se encuentra este proyecto. Se describe su arquitectura general y los tipos de módulos que lo componen, los drones utilizados, y se explica la contribución de este trabajo al bloque de aplicaciones Flutter del ecosistema.

El Capítulo 3 agrupa el marco tecnológico del proyecto. Presenta el framework Flutter y el lenguaje de programación Dart, los modos de comunicación del DEE, las interfaces de programación de aplicaciones (API), y los protocolos de transporte y de aplicación usados (TCP/UDP, HTTP/HTTPS, MQTT y WebRTC). El capítulo concluye con los repositorios y herramientas de desarrollo utilizados a lo largo del proyecto,

El Capítulo 4 recoge los objetivos del proyecto y el plan de trabajo. Se detalla el objetivo general, los objetivos específicos y una descripción ordenada de las tareas principales a realizar.

El Capítulo 5 presenta la arquitectura software base capaz del control mínimo del dron. 

El capítulo 6 documenta cada una de las demostraciones de paradigmas de control desarrolladas de forma independiente junto a su implementación en EZDrone y una demostración de su funcionamiento. 

El Capítulo 7 describe la evolución de la arquitectura de comunicación, desde el diseño inicial (Capítulo 5) hasta el diseño híbrido definitivo. 

El Capítulo 8 presenta EZDrone en su conjunto, detallando la arquitectura, los detalles de implementación y los resultados de las pruebas. 

El Capítulo 9 recoge el estudio de la sostenibilidad. 

El Capítulo 10 cierra el documento con las conclusiones y las líneas de trabajo futuro.


\chapter{Introducción al DEE}   % capítulos numerados normales
Este capítulo presenta el Drone Engineering Ecosystem (DEE), la plataforma académica desarrollada en la EETAC. El apartado 2.1 describe la arquitectura general del ecosistema y los tipos de módulos que lo componen. El apartado 2.2 presenta los drones utilizados en el DEE. El apartado 2.3 explica dónde se sitúa la contribución de este proyecto dentro del ecosistema.

\section{Arquitectura del DEE}

El DEE es una plataforma software abierta y modular que permite el control y la monitorización de uno o varios drones mediante diferentes tipos de dispositivos y aplicaciones. Su arquitectura se organiza en torno a cuatro tipos de módulos.

Los módulos a bordo se ejecutan en una Raspberry Pi montada sobre la plataforma del dron y se encargan de la interfaz directa con el autopilot, la cámara, los LEDs y otros periféricos. Para facilitar el despliegue, estos módulos se empaquetan en contenedores Docker.

Las aplicaciones frontend proporcionan la interfaz al operador e incluyen aplicaciones de escritorio (Python/TKinter, C\#), aplicaciones web (Flask, Vue.js) y aplicaciones móviles nativas (Android Studio, Flutter).

Los módulos backend se ocupan del almacenamiento y recuperación de datos.

La infraestructura de comunicación está formada por un broker interno y uno externo (público o alojado en la UPC), que conecta todos los módulos a través de internet. El modo de comunicación principal para las aplicaciones web y el despliegue remoto es el modo global, que enruta todos los mensajes a través del broker externo. El DEE soporta también un modo local, cuando no hay cobertura de internet, y un modo directo, empleado en simulación Software In The Loop (SITL), donde todos los módulos se ejecutan en un único dispositivo.


\section{Drones utilizados en el DEE}

El dron principal usado en el DEE consiste en un chasis HexSon con motores brushless y un controlador de vuelo Orange Cube PixHawk ejecutando firmware ArduPilot\cite{hexsoon}. La plataforma incorpora además una cámara. Este dron está destinado a operaciones en exteriores y es la plataforma de referencia para el control de todos los proyectos como EZDrone.

El DJI Tello es un dron pequeño, destinado a misiones en interiores. Dentro del DEE se utiliza principalmente en demostraciones a visitantes, dado que su tamaño reducido y su comportamiento estable lo hacen especialmente adecuado para entornos cerrados\cite{telloEETAC}. EZDrone incluye soporte para el control del Tello, permitiendo realizar estas demostraciones directamente desde el navegador.


\section{Contribución de este proyecto al DEE}
Este proyecto contribuye al bloque de aplicaciones Flutter del DEE con dos aportaciones complementarias. La primera es una colección de aplicaciones Flutter de demostración. Cada demostración está documentada de forma autónoma y está pensada para futuros desarrolladores. La segunda aportación (EZDrone), cae en webs de Flutter para el control de drones.


\chapter{Flutter y fundamentos}
Este capítulo presenta el marco tecnológico del proyecto. El apartado 3.1 introduce el framework de Flutter y el lenguaje de programación Dart. El 3.2 explica los modos de comunicación del DEE. El 3.3 introduce el concepto de API y los tipos más relevantes. Los apartados 3.4-3.7 describen los protocolos de transporte y de aplicación usados: TCP/UDP, HTTP/HTTPS, MQTT y WebRTC. El capítulo concluye con el apartado 3.8, que agrupa los repositorios y herramientas de desarrollo utilizados a lo largo del proyecto.

\section{Framework Flutter}
Flutter es un framework de código abierto desarrollado por Google para la construcción de interfaces de usuario nativas y compiladas desde una única base de código. A diferencia de los frameworks tradicionales, que generan una interfaz web envuelta en un contenedor, Flutter dispone de su propio motor de renderizado, y dibuja cada píxel directamente en el canvas del sistema operativo, sin tener que delegar en los componentes de cada plataforma. Desde una misma base de código, es posible compilar y desplegar la aplicación en Android, iOS, web, escritorio (Windows, macOS, Linux) y dispositivos embebidos.

\subsection{Lenguaje Dart}
Flutter utiliza Dart como lenguaje de programación\cite{dart}. Dart es un lenguaje compilado, orientado a objetos, también desarrollado por Google, con soporte nativo para programación asíncrona, mediante los comandos \texttt{\'"Async / Await"} y la clase \texttt{"Future"}, y para flujos de datos continuos mediante \texttt{"Stream"}. Para la compilación web, el compilador  \texttt{"dart2js"}\cite{dart2js} traduce el código Dart a JavaScript optimizado, que es lo que se ejecuta en el navegador en el caso de Flutter Web.

\subsection{Modelo de widgets y árbol de renderizado}
La unidad fundamental de construcción en Flutter es el widget\cite{flutter_arch}. Todo en la interfaz es un widget: los botones, el texto, el mapa, los contenedores del layout e incluso los gestores de estado. Los widget son objetos inmutables que describen una porción de la interfaz, no la representan directamente, y Flutter reconstruye solo las partes del árbol que han cambiado cuando el estado de la aplicación se modifica.

Flutter distingue dos tipos de widgets. Los \texttt{StatelessWidget}, widgets sin estado interno: su apariencia depende únicamente de los parámetros que reciben y del estado del sistema proporcionado por sus ancestros en el árbol. Son los más ligeros y se utilizan para componentes cuya representación no cambia una vez construidos. Los \texttt{StatefulWidget} contienen un objeto State mutable: cuando se llama a \texttt{setState( )}, Flutter invalida el widget y lo reconstruye con los nuevos valores. En este proyecto, la pantalla de configuración y la de vuelo de la arquitectura base son \texttt{StatelessWidget} suscritos al \texttt{DronProvider}, mientras que el botón de movimiento continuo \texttt{ContinuousButton} es un \texttt{StatefulWidget} porque necesita rastrear si está siendo pulsado en cada instante.

El proceso de renderizado de Flutter se divide en tres árboles que se mantienen sincronizados: el árbol de widgets (la descripción declarativa de la interfaz), el árbol de elementos (la instancia persistente de cada widget en el árbol) y el árbol de renderizado (los objetos que calculan el layout, la pintura y la composición). Este diseño permite que las reconstrucciones del árbol de widgets no sean muy pesadas, ya que Flutter reutiliza los elementos y objetos siempre que no hayan cambiado.

\subsection{Gestión de estados con Provider}
Para aplicaciones que deben reflejar un estado compartido entre múltiples widgets, como la telemetría del dron, el mapa y los botones de la barra inferior, Flutter no impone un patrón específico, pero la comunidad de usuarios de Flutter ha convergido en varias soluciones. Este proyecto utiliza la librería Provider\cite{provider}, que implementa el patrón de inyección de dependencias y notificación de cambios sobre el mecanismo nativo \texttt{InheritedWidget} de Flutter.

La clase \texttt{DronProvider} centraliza todo el estado de la aplicación, las variables de telemetría y los flags de estado, los métodos que ejecutan las acciones del operador y la lógica de obtención de datos del dron. Cuando cualquiera de estas variables cambia, la llamada a \texttt{notifyListeners()} provoca la reconstrucción de todos los widgets que se han suscrito al provider mediante \texttt{context.watch<DronProvider>()}. Los widgets que solo necesitan usar métodos del provider, sin suscribirse a sus cambios, utilizan \texttt{context.read<DronProvider>( )}, que no establece una suscripción. Es el acceso correcto dentro de callbacks como el \texttt{onPressed} de un botón.

Usando este diseño, las pantallas no contienen ninguna lógica de estado y se limitan a leer el Provider y a invocar sus métodos. 

\section{Modos de comunicación}
El DEE define tres modos de comunicación que determinan cómo se conectan entre sí los módulos a bordo, las aplicaciones de frontend y los servicios de backend.

El modo global asume que todos los componentes están conectados a internet. La comunicación se enruta a través de un broker MQTT externo (HiveMQ\cite{hivemq} o el broker de la UPC). Es el modo estándar para la operación remota y para aplicaciones web, y el que emplea EZDrone.

El modo local se utiliza cuando no hay cobertura de internet. El dispositivo del operador se conecta directamente al punto de acceso Wi-Fi proporcionado por la Raspberry Pi a bordo, y una instancia de Mosquitto\cite{mosquitto} se ejecuta localmente. Los módulos del frontend se conectan a 10.10.1.8000 y los módulos a bordo a localhost:8000.

El modo directo se emplea cuando no existe computador a bordo. El servicio de autopilot se ejecuta como parte de la aplicación de frontend y se conecta directamente al autopilot mediante una radio de telemetría. Es también el modo utilizado en la simulación SITL\cite{ardupilot_sitl}, donde todos los módulos, incluyendo el broker, se ejecutan en un único dispositivo.

\section{Interfaces de programación de aplicaciones (API)}
Una API (Application Programming Interface)\cite{api} es un contrato software que define cómo dos componentes se comunican entre sí, qué operaciones están disponibles, qué formato tienen los mensajes y qué respuestas pueden esperarse. En el contexto de este proyecto, las APIs relevantes son de tipo red, es decir, definen el intercambio de mensajes entre procesos que se ejecutan en dispositivos distintos.

Los dos tipos de APIs de red más utilizados en el desarrollo web son las APIs REST y las APIs basadas en WebSocket. EZDrone usa ambas en distintas fases de su evolución arquitectónica.

Una API REST (Representational State Transfer)\cite{apirest} es un estilo arquitectónico que opera sobre HTTP. Cada recurso o acción del sistema se expone como una URL, y las operaciones se expresan mediante verbos HTTP estándar: \texttt{GET} para consultar un estado, \texttt{POST} para enviar comandos. La interacción es siempre iniciada por el cliente, y es sin estado, el servidor no retiene información entre peticiones. En la arquitectura base descrita en el Capítulo 5, el backend Dart expone una API REST sobre HTTP con endpoints como \texttt{/connect}, \texttt{/arm}, \texttt{/takeoff} o \texttt{/status}. El frontend Flutter llama periódicamente a \texttt{/status} para obtener telemetría.

WebSocket\cite{websocket} es un protocolo que establece un canal de comunicación full-duplex y persistente sobre una única conexión TCP. A diferencia de HTTP, una vez completado el handshake inicial, tanto el cliente como el servidor pueden enviar mensajes en cualquier momento y en ambas direcciones, sin que el otro extremo haya realizado una petición previa. Esto elimina la necesidad de estar constantemente solicitando actualizaciones de estado y reduce la latencia a la del propio canal. 

\section{Protocolos de transporte: TCP y UDP}
Antes de describir los protocolos de aplicación empleados en este proyecto, es necesario presentar los dos protocolos de capa de transporte sobre los que se construyen, TCP y UDP, cuyas diferencias fundamentales justifican las decisiones arquitectónicas adoptadas en el proyecto.

TCP (Transmission Control Protocol)\cite{tcp} es un protocolo orientado a la conexión que garantiza una entrega fiable y ordenada de los datos entre dos dispositivos. Antes de intercambiar ningún dato, se establece una conexión mediante un proceso de handshake. A partir de ese momento, todos los paquetes enviados son confirmados por el receptor, si esta confirmación no llega, el paquete se retransmite automáticamente. Este mecanismo garantiza que ningún dato se pierde, duplica o entrega fuera de orden. El coste de esta fiabilidad es una alta latencia (culpa de los ciclos de confirmación y retransmisión, ya que si hay una retransmisión, todos los paquetes posteriores quedan retenidos).

UDP (User Datagram Protocol)\cite{udp} adopta el enfoque contrario, es un protocolo sin conexión que envía paquetes sin establecer una conexión previa ni ofrecer ninguna garantía de entrega. Los paquetes pueden perderse, llegar fuera de orden o duplicarse, y el protocolo no realiza ningún intento de corregirlo. Lejos de ser una limitación, esta característica hace que sea idóneo para flujos continuos en tiempo real. Si un paquete que transporta una trama del joystick o un segmento de vídeo se pierde, lo más eficiente es descartarlo y utilizar el siguiente paquete, una lectura ligeramente desactualizada es menos perjudicial que un flujo paralizado esperando una retransmisión.

Este balance entre fiabilidad y latencia, hace que en el proyecto se adopte el uso de los dos protocolos: MQTT, construido sobre TCP, se emplea para los comandos discretos y las notificaciones de cambio de estado, donde la pérdida de un mensaje es inaceptable. WebRTC construido sobre UDP, se emplea para la telemetría continua y el vídeo, donde la latencia importa más que la entrega garantizada.

\section{Protocolo HTTP}
HTTP (Hypertext Transfer Protocol)\cite{http} es el protocolo de comunicación de datos de la web. Opera sobre TCP y sigue un modelo de petición-respuesta. El cliente envía una petición, el servidor la procesa y devuelve una respuesta, justo después, la conexión se cierra o se mantiene para recibir más peticiones. HTTP es un protocolo sin estado, lo que significa que cada petición es independiente y el servidor no tiene memoria de las interacciones anteriores.

En el contexto de este proyecto, HTTP es el protocolo empleado en la arquitectura base descrita en el Capítulo 5. El frontend Flutter emite peticiones \texttt{HTTP POST} para lanzar comandos y peticiones \texttt{HTTP GET} para consultar el estado del dron. Esta última estrategia recibe el nombre de "polling", consultar periódicamente el servidor en busca de datos actualizados. Es sencilla de implementar y suficiente para actualizaciones de baja latencia, pero introduce latencia proporcional al intervalo de consulta y genera tráfico de red innecesario cuando no hay cambio de estado. Esta limitación motiva directamente la evolución arquitectónica descrita en el Capítulo 7, donde el polling HTTP es reemplazado por canales de datos WebRTC.

\subsection{HTTPS y TLS}
HTTPS es la versión segura y certificada del HTTP. No se trata de un protocolo distinto, sino que es el mismo HTTP con una capa de seguridad añadida mediante TLS (Transport Layer Security)\cite{tls}, que cifra todo el tráfico entre el cliente y el servidor. Cualquier dato intercambiado es ilegible para un tercero que intercepte la comunicación, ya que solamente los dos extremos de la conexión tienen las claves necesarias para descifrarlo.

El establecimiento de una conexión TLS se realiza mediante un handshake, donde el servidor presenta su certificado digital (emitido por una autoridad certificadora, "CA", de confianza que acredita que el servidor es quien dice ser y que contiene su clave pública). El cliente verifica que el certificado haya sido acreditado por una CA, y a partir de allí ambos extremos negocian una clave de sesión con la que cifrarán el resto del intercambio. 

En el contexto de este proyecto, el uso de HTTPS es obligatorio, ya que los navegadores bloquean muchas acciones hechas en HTTP plano, ya sea el uso de WebSockets seguros, o APIs con acceso interno a los dispositivos (como el control por IMU o por voz)\cite{bloqueo_http}.

\section{Protocolo MQTT}
MQTT (Message Queuing Telemetry Transport)\cite{mqtt} es el protocolo principal de intercambio de comandos y estados entre módulos en el DEE. Se trata de un protocolo de mensajería publicación-suscripción ligero.

El protocolo opera en torno a un broker central al que se conectan todos los módulos. Un módulo que desea enviar un mensaje lo publica en un topic específico, cualquier módulo suscrito a ese topic recibe el mensaje. Este modelo desacoplado implica que publicadores y suscriptores no necesitan conocerse mutuamente, únicamente comparten el broker y la convención de nombre de los topics.

En el DEE, los nombre de topics siguen un convenio estandarizado con la siguiente forma:

\texttt{[módulo\_origen] / [módulo\_destino] / [comando]}

Por ejemplo, cuando la aplicación Flutter solicita un despegue, publica en el topic \texttt{"mobileFlutter/ groundStation/ takeoff"}. La estación de tierra, suscrita a \texttt{"mobileFlutter/ groundStation/  \#"}, recibe el mensaje y emite la orden correspondiente al autopiloto. Cuando el dron ha alcanzado la altitud objetivo, la estación tierra publica en \texttt{groundStation/ mobileFlutter/ flying}. Esta convención bidireccional hace el enrutamiento de mensajes explícito y trazable, lo que resulta útil durante el desarrollo y la depuración.

Para las aplicaciones Flutter Web, la comunicación MQTT debe utilizar el transporte WebSocket en lugar de TCP directo, dado que los navegadores no permiten abrir conexiones de socket arbitrarias. El broker público de HiveMQ soporta ambas variantes del WebSocket, WebSocket plano (puerto 8000) y WebSocket Seguro (WSS, puerto 8884), siendo este último obligatorio cuando la aplicación web se sirve mediante HTTPS.

\section{Protocolo WebRTC}
WebRTC (Web Real-Time Communication)\cite{webrtc} es un estándar abierto y una API que permite la comunicación peer-to-peer de audio, vídeo y datos directamente entre navegadores y aplicaciones nativas, sin requerir de servidores intermedios para los flujos de datos en sí. Desarrollado originalmente para la videoconferencia en navegadores, ha sido adoptado ampliamente en teleoperación, robótica, streaming en vivo y muchas aplicaciones, gracias a su baja latencia y su soporte nativo en todos los navegadores modernos.

Una conexión WebRTC se establece mediante un proceso de señalización (o signalling)\cite{signalling} en el que los dos peers intercambias descriptores de sesión (ofertas y respuestas SDP, Session Description Protocol) e informan de accesibilidad de red (candidatos ICE, interactive Connectivity Establishment) a través de un servidor de señalización intermediario . Una vez establecida la conexión, los flujos medios y datos circulan directamente entre peers. El servidor de señalización sólo es necesario para la fase del handshake.

En EZDrone, el servidor de señalización es el mismo backend Dart, alojado en el dominio de la EETAC \texttt{dronseetac.upc.edu:8105}. Este servidor, además de retransmitir los mensajes SDP e ICE entre la estación tierra Python y el cliente Flutter, envía a cada peer recién conectado la configuración ICE del ecosistema: un servidor STUN público y un servidor TURN privado alojado en dronseetac.upc.edu, necesario para establecer la conexión cuando alguno de los peers se encuentra detrás de un NAT estricto.

En este proyecto, se usan dos tipos de canales WebRTC,

\texttt{Las pistas de vídeo/audio, son flujos de medios unidireccionales para la transmisión de la imagen de la cámara desde }ARRREEEEGGGGGLLLLLAAAARRRRR

Los DataChannels son canales de datos bidireccionales y de baja latencia, para la transmisión de datos binarios o de texto arbitrarios entre peers. En EZDrone se utilizan dos Datachannels, el de los joysticks, por el que Flutter envía los valores de los cuatro ejes del mando a la estación tierra, y el de la telemetría, por el cual la estación tierra envía a Flutter el estado del dron y los datos de vuelo en formato JSON\cite{json}.

La ventaja principal de utilizar DataChannels WebRTC para la telemetría y el control de joysticks frente a MQTT reside en la latencia y el desacoplamiento del broker, una vez se ha establecido la conexión entre los peers, los datos fluyen directamente entre ellos sin pasar por el broker, eliminando el retardo de ida y vuelta, y el riesgo de saturación del broker ante actualizaciones de tan alta frecuencia.

\section{Repositorios y herramientas de desarrollo}
El DEE se organiza en repositorios GitHub, uno por bloque funcional. Cada repositorio contiene el código fuente, instrucciones de instalación, tutoriales y vídeos de demostración. El repositorio principal, que sirve como punto de entrada y proporciona la descripción de la arquitectura global, es https://github.com/dronsEETAC/DroneEngineeringEcosystemDEE. 
El bloque directamente relevante para este proyecto es el repositorio de aplicaciones Flutter, que consolida todos los módulos de frontend basados en Flutter contribuidos por estudiantes en sucesivos proyectos de fin de estudios.

Las herramientas de desarrollo principales empleadas a lo largo del proyecto son: VS Code con los plugins de Flutter y Dart (y muchos más explicados a lo largo del proyecto) para el desarrollo del frontend, el backend de señalización y la estación de tierra en Python junto a la librería de DronLink. Mission Planner en Windows para la simulación SITL. El broker público de HiveMQ como broker MQTT externo y Git y GitHub para el control de versiones. 

\chapter{Objetivos, tareas y plan de trabajo}

\section{Objetivos}
\subsection{Objetivo general}
El objetivo principal de este proyecto es diseñar, implementar y documentar una colección de aplicaciones Flutter de demostración y una estación de control de tierra web (EZDrone), destinadas al control de drones ArduPilot y DJI Tello dentro del DEE. 

\subsection{Objetivos específicos}
\begin{itemize}
    \item {OE1- Control básico de un dron.}Implementar una aplicación Flutter base capaz de conectarse a un dron Ardupilot mediante MQTT y ejecutar las operaciones de vuelo fundamentales: armado de motores, despegue, navegación direccional, retorno al punto de lanzamiento (RTL) y aterrizaje.
    \item {OE2 - Demostraciones de paradigmas de control.}Desarrollar demostraciones Flutter independientes y documentadas que exploren los distintos mecanismos de control  e interacción con el dron: joystick virtuales, control por voz o mediante mediciones inerciales de la IMU, localización GPS del operador, streaming de vídeo con detección de objetos, planificación de misiones autónomas y registro y análisis de vuelos.
    \item {OE3 - Arquitectura de comunicación híbrida.}Diseñar una arquitectura que combine MQTT (para comandos discretos y cambios de estado) con canales de datos y pistas de vídeo WebRTC (para parámetros de alta frecuencia como telemetría continua y vídeo en tiempo real).
    \item {OE4 - Estación de tierra en Python.}Implementar una estación de tierra en Python que traduzca los comandos normalizados recibidos vía MQTT hacia las llamadas correspondientes de la biblioteca DroneLink, actuando como puente entre la interfaz Flutter y el autopiloto. 
    \item {OE5 - Validación mediante SITL.}Validar el sistema completo mediante el simulador ArduPilot SITL (software In The Loop), confirmando la correcta ejecución de todos los comandos de vuelo y el control en tiempo real, antes de las pruebas con dron real.
    \item {OE6 - Repositorio abierto a todos.}
    Publicar un repositorio estructurado y públicamente accesible que contenga todas las demostraciones y el proyecto EZDrone final, organizado como recurso de referencia práctica para la comunidad.
\end{itemize}

\section{Plan de trabajo}
El desarrollo de este proyecto se ha estructurado en una secuencia de tareas principales. A continuación se describe brevemente cada una de estas fases.
\begin{itemize}
    \item {Tarea 1 - Aplicación de control básico.}Implementación de la primera versión funcional del sistema: frontend Flutter backend Dart y estación tierra Python. Esta fase establece la arquitectura base, el modelo de gestión de estados con Provider y la cadena de comandos HTTP - MQTT - DronLink. La validación se realiza mediante SITL.
    \item {Tarea 2 - Desarrollo de la colección de demostraciones.}Desarrollo de las demostraciones Flutter de forma secuencial e independiente, joystick virtuales, control por voz, control inercial mediante IMU, localización GPS en mapa, streaming de vídeo WebRTC con detección YOLO, planificación de misiones autónomas, registro y análisis de vuelos y finalmente catálogo de widgets (éste último se ha realizado al finalizar EZDrone). Cada demostración se documenta de forma autónoma e incluye instrucciones de integración.
    \item {Tarea 3 - Evolución de la arquitectura de comunicación.}Eliminación del backend Dart como intermediario y conexión y comunicación. Incorporación de WebRTC para los canales de alta frecuencia (joysticks, telemetría, vídeo). Implementación y despliegue del servidor en servidores de la EETAC (dronseetac.upc.edu).
    \item {Tarea 4 - Integración de las demostraciones en EZDrone.}
    Integración de todos los paradigmas de las demostraciones en la web EZDrone, incluyendo también el soporte para los drones ArduPilot y Tello, no solo SITL. Implementación de las pantallas de vuelo en sus diferentes modos, planificación de misiones e historial de vuelos.
    \item {Tarea 5 - Validación con SITL y dron real.}Validación del sistema completo mediante simulación SITL y pruebas reales con los dos drones, verificando el correcto comportamiento de todos los modos de control y la robustez de la comunicación.
    \item {Tarea 6 - Redacción de la memoria y publicación de los repositorios.}Redacción de la memoria y estructuración y publicación del repositorio de demostraciones y del proyecto EZDrone, con documentación e instrucciones para futuros desarrolladores del ecosistema.
\end{itemize}
\chapter{Aplicación de control básico}
Este capítulo presenta el primer hito funcional del proyecto, una aplicación Flutter Web capaz de conectarse al dron Ardupilot, armar sus motores, despegar, navegar en las seis direcciones básicas y aterrizar. Esta versión establece la arquitectura base y el modelo de gestión de estados sobre los que se basarán todas las arquitecturas posteriores. El apartado 5.1 describe la arquitectura general del sistema. El 5.2 detalla la cadena de comandos seguida desde la interfaz del usuario hasta el dron. El 5.3 explica la máquina de estado del dron y su gestión mediante un Provider. El 5.4 presenta las pantallas de la interfaz. El último apartado, 5.5, describe el proceso de validación mediante SITL.

\section{Arquitectura del sistema}
La primera versión está compuesta por tres componentes que se comunican entre sí mediante dos canales independientes. El frontend Flutter, el bakcend Dart y la estación tierra en Python.

El frontend Flutter Web es una aplicación que se ejecuta en el navegador y se estructura en torno a un gestor de estados centralizados, la clase \texttt{Provider}. Se comunica con el backend Dart mediante peticiones HTTP, peticiones \texttt{POST} para emitir comandos de vuelo y peticiones \texttt{GET} periódicas para obtener el estado actual del dron y los datos de telemetría (el polling mencionado anteriormente en la sección 3.5).
El backend Dart es un servidor HTTP construido con el framework Shelf, que escucha en el puerto 8080 y contiene una API REST. Recibe los comandos del frontend, los traduce en publicaciones MQTT hacia el broker HiveMQ, y mantiene actualizado al frontend sobre el estado del dron (frontend lo consulta mediante el polling).

La estación tierra Python es un script que se conecta al mismo broker que el servidor (broker.hivemq.com, puerto 1883), se suscribe al árbol de topics, interpreta los comandos recibidos y emite las instrucciones correspondientes al dron mediante la librería DronLink. Además de eso, publica confirmaciones de estado, y periódicamente la telemetría de vuela al servidor.
Esta arquitectura de tres componentes con el backend como intermediario es coherente con el modo de comunicación Global del DEE, descrito en el Capítulo 2. 

\section{Cadena de comandos}
Los comandos usados se mandan desde la interfaz Flutter hasta el dron a través de una pipeline de dos tramos. El primero, se trata de un tramo HTTP que comunica el frontend y el backend, el segundo, es un tramo MQTT que junta el backend con la estación tierra.

Cuando el usuario pulsa un botón en la interfaz, el DronProvider invoca el método correspondiente de la clase HttpLogic. Esta clase centraliza todas las llamadas HTTP y usa la librería http de Dart. Para los comandos sin parámetros (como son Connect, Arm, Land, Rtl o Disconnect), se utiliza el método privado sendPostRequest(url), que emite un POST y verifica que la respuesta sea 200 (OK). Para los comandos con parámetro (como takeoff o speed), se serializa el valor en la payload de la petición como JSON. Finalmente, para recibir, se usa el método getStatus( ), que realiza un GET al endpoint /status y devuelve una payload con los datos.

El backend, define un router Shelf con un handler. Al recibir una petición, cada handler actualiza el estado interno según corresponda, y publica el topic MQTT equivalente hacia la estación tierra. Para poner un ejemplo, cuando se envía el comando de armar, el handler de /arm, pone los flags de dronFlying y dronArmed en false como medida de precaución antes de publicar el /arm. O por ejemplo, el handler de /takeoff lee el campo altitude de la payload JSON y lo publica en /takeoff. 

Todos los handlers devuelven un código 200 (junto a las cabeceras CORS necesarias para permitir las peticiones desde el navegador) al frontend.

La estación tierra, suscrita a mobileFlutter/groundStation/\#, recibe todos los mensajes que publica el backend, y llama a los métodos de Dronlink correspondientes (como podría ser /arm - dron.arm( ), /takeoff - dron.tafeOff(altitude) o /land - dron.Land( ) ). Una vez la operación se ha realizado, la estación tierra publica la confirmación en el topic de respuesta (groundStation/mobileFlutter/armed,flying,landed…).

El camino de retorno de la telemetría sigue el modelo de polling. La estación tierra llama periódicamente a dron.send\_telemetry\_info( ), que construye un diccionario con todos los datos de telemetría, lo serializa en formato JSON y lo publica en groundStation/mobileFlutter/telemetry, donde el backend los recibirá y almacenará. Finalmente, en el siguiente ciclo de polling, el frontend obtendrá dichos valores mediante un GET a /status.

\section{Máquina de estados}
El estado del dron se rastrea de forma independiente en los tres componentes de la aplicación. La fuente de mayor autoridad es la estación tierra, ya que tiene acceso directo al autopiloto del dron. El backend mantiene una réplica de este, actualizado mediantes los topics MQTT que envía la estación tierra. El frontend tiene su propia representación de los estados mediante el DronProvider.

Las variables de estado que gestiona el Provider son principalmente tres booleanos: isConnected, que pasa a verdadero cuando la estación tierra confirma la conexión con el dron, isArmed, que pasa a verdadero cuando se ha confirmado que los motores están armados y listos para el despegue, e isFlying, que permanece en verdadero desde la confirmación del despegue hasta la confirmación del aterrizaje. Estas variables son las que tienen el control de qué botones están disponibles en la interfaz. Por ejemplo, “Takeoff” solo está habilitado cuando isConnected e isArmed són verdaderos, o “Land” y “RTL” solo están habilitados cuando isFlying es true.

Un reto que ha sido destacable, es la confirmación del armado. ArduPilot ejecuta una secuencia de comprobaciones previas al armado, que pueden tardar varios segundos en completarse. El DronProvider se encarga de gestionar esta situación de la siguiente manera: Se implementa una deadline, cuando se pulsa ARM se activa un boolean “-waitingForArm” y empieza un contador de 5 segundos en “-armDeadline”. Si en uno de los ciclos de polling se detecta que armed = true, entonces se cancela la espera, pero si el plazo se termina sin la confirmación del armado, se entiende que no se ha podido armar y se muestra el mensaje “Not ready to arm”. Este mecanismo evita que la interfaz se desarme automáticamente o se quede en un estado de carga indefinido.

La telemetría recibida en cada ciclo de polling se almacena en una serie de variables de tipo double y String. Adicionalmente, cada posición GPS válida (distinta a 0 o null), se guarda en una lista “droneTrail”, para así representar la trayectoria del dron en el mapa.

\section{Interfaz de usuario}
La aplicación se compone de dos pantallas con navegación secuencial, implementadas como widgets sin estado (StatelessWidget) suscritos al DronProvider mediante “context.Watch<DronProvider>( ).

\subsection{Setup Screen}
La pantalla de configuración (SetupScreen) es el punto de partida de la aplicación. Contiene dos campos de texto para la configuración del vuelo, la altitud de despegue (con un rango válido entre 2,0 y 50,0 metros) y la velocidad de vuelo (con rango entre 1,0 y 15,0 metros por segundo). Ambos campos se validan en tiempo real gracias al Provider, que les atribuye una flag “isConfigValid” donde si se introduce un valor dentro de los márgenes se iguala a True. Si el valor está fuera del rango, el texto auxiliar se muestra en rojo, y el botón “START FLIGHT” permanece deshabilitado hasta que se introduzcan valores válidos.

La pantalla incluye tres botones de gestión: 
\begin{itemize}
  \item \textbf{Connect:} Llama a connectDron( ), que envía un POST a /connect, pone al provider en “isLoading” = True e inicia el polling.
  
\item \textbf{Start flight:} Navega a la pantalla de vuelo (mediante un Navigator.push, herramienta navegadora de Dart), y sólo está habilitado cuando isConnected e isConfigValid son True.

\item \textbf{Disconnect:}  Llama a diconnectDron( ), que envía un POST a /disconnect, detiene el polling y limpia todas las variables.
\end{itemize}

\subsection{Flight Screen}
La pantalla de vuelo (Flight Screen) es la interfaz operativa. Se organiza en zonas horizontales (rows).

La barra superior muestra la telemetría en tiempo real: altitud (ALT), velocidad (GS), rumbo magnético (HDG), latitud (LAT), longitud (LON) con 7 decimales (igual que en Mission Planner), estado del autopiloto (STATE), modo de vuelo (MODE) y nivel de batería. Este último, tiene un código de colores, donde el símbolo de la batería aparece verde por encima del 50%, naranja entre el 20% y 50% y rojo por debajo del 20%. En el centro de esta zona se encuentra un texto que muestra el “message” del Provider.

La zona central presenta un mapa satelital proporcionado por la librería flutter\_map con tiles de OpenStreetMap. El mapa muestra el icono del dron en su posición GPS actual, rotado según su rumbo “currentHeading”. Una flecha naranja superpuesta indica la dirección de vuelo, y la trazada del dron se indica con una polilínea púrpura “droneTrail”. La velocidad vectorial la indica una línea verde formada por el vector entre los campos de velocidad Vx y Vy, que aparte de mostrar la magnitud relativa de la velocidad, también indica la dirección del movimiento actual. Y finalmente una sombra circular girs que indica si el dron está volando.

Si el dispositivo del operador tiene señal GPS, se muestra su posición como un icono de móvil azul junto con un indicador de precisión codificado en colores (verde por debajo de 5 m, naranja por debajo de 20 m, rojo en caso superior).

A ambos lados del mapa se sitúan los controles de movimiento, implementados como un widget de ContinuousButton. Este widget utiliza la función propia de Dart GestureDetector para detectar onTapDown y onTapUp, de forma que mientras el botón esté pulsado, va llamando a DronProvider.startmove(direction), que envía un POST a /move con la payload JSON {direction: “Forward”} (por ejemplo). Al soltar el ContinuousButton, se llama a DronProvider.stopMove( ), que envía la dirección “Stop”, parando el movimiento del dron.
El widget tiene una animación al pulsarlo (usando AnimatedContainer), donde se reduce el tamaño y se cambia el color, para así indicar visualmente cuál se está pulsando. 
En esta primera versión, los controles de movimiento son flechas. En futuras versiones serán cambiadas por joysticks virtuales.

La barra inferior contiene cinco botones de vuelo principales: ARM, TAKEOFF, DISCONNECT, LAND y RTL. La habilitación de cada botón viene determinada por la combinación de “isLoading”, “isConnected”, “isArmed” e “isFlying”, garantizando que solo son accesibles los comandos válidos para cada estado del dron (es decir, que si está volando, “isFlying” = True, no se pueda desonnectar, o no pueda armarse si ya está armado, “isArmed” = True).

\section{Validación mediante SITL}

La aplicación base fue validada usando Mission Planner como emulador de vuelo en modo SITL (Software In The Loop). Esta modalidad permite ejecutar el firmware ArduPilot en el propio ordenador de desarrollo, simulando un vehículo con comportamiento realista sin necesidad de usar un dron real. La estación tierra Python se conecta al simulador mediante la cadena de conexión “tcp:127.0.0.1: 5763", que es el puerto TCP estándar que expone el SITL del ArduPilot.

La validación confirmó la correcta ejecución de la secuencia de vuelo completa: establecimiento de la conexión, armado de motores, despegue a la altitud configurada, navegación en las seis direcciones, retorno al punto de lanzamiento y aterrizaje. También permitió verificar el correcto funcionamiento del mecanismo de deadline del armado y la actualización continua de la telemetría en la interfaz durante todo el vuelo.

\chapter{Demostraciones de paradigmas de control}
Este capítulo presenta la colección de aplicaciones Flutter de demostración desarrolladas como contribución de referencia del proyecto al DEE. Cada demostración explora un paradigma específico de control o monitorización del dron de forma autónoma. El objetivo es proporcionar un ejemplo mínimo, bien documentado y reutilizable para cada tecnología, que pueda ser estudiado, adaptado e integrado de forma independiente por futuros desarrolladores del ecosistema.

El apartado 4.1 describe la organización del repositorio. Los apartados 4.2 a 4.X presentan cada demostración individualmente: joysticks virtuales, control por voz, control inercial mediante IMU, localización GPS en mapa y streaming de vídeo WebRTC con detección de objetos.

\section{Organización y estructura del repositorio}
Todas las demostraciones se organizan como proyectos Flutter independientes, dentro de la carpeta “demostraciones/” del repositorio del proyecto. Cada proyecto dispone de su propio “pubspec.yaml” y de su propio “main.dart” como punto de partida. Esta independencia es un requisito de diseño elegido, ya que permite que un desarrollador que solo necesite entender el control por voz, por ejemplo, pueda clonar el proyecto, abrir la carpeta que toca y ejecutar la demostración sin necesidad de configurar el resto del sistema.

Todos los ficheros de código están comentados, documentando no solo el funcionamiento sino también como debería adaptarse para conectarse a un dron real. Este enfoque orientado a la reutilización es lo que distingue esta colección de un conjunto de prototipos, y la convierte en un recurso de referencia para la comunidad.

\section{Catálogo de widgets}

\subsection{Caso de uso}
Una estación de control usa una paleta reducida, pero muy específica de componentes de interfaz. Botones que cambian de color según el estado, controls que responden a un pulsado contínuo, sliders de zoom, selectores de modo, listas reordenables etc. Cada uno de estos widgets tiene un comportamiento, unos estados visuales y unas características que no són evidentes a primera vista y que, si se implementan incorrectamente pueden causar varios problemas. 

Esta demostración es un poco distinta a las que le proceden. No explora un paradigma de control sino que actúa como catálogo interactivo y ejecutable de los 10 tipos de widget principales empleados en EZDrone, con cada igual a como se encuentran en este, con sus variantes visuales documentadas y una nota que explica los detalles de implementación relevantes. El objetivo es que un desarrollador que quiera reutilizar algún componente pueda ver el widget en acción, entender su lógica y copiarlo.

\subsection{Implementación}
La demostración se estructura como una pantalla de scroll vertical con diez secciones, cada una encabezada por el nombre del widget, su uso en EZDrone y una nota que documenta los aspectos no evidentes de su implementación.

El primer widget es ElevatedButton.icon, el botón de acción principal usado en todas las pantallas de EZDrone. La sección replica la barra de controles de vuelo (arm, take off, land, rtl y disconnect). Cada botón tiene habilitación condicional: arm solo está disponible cuando isConnected es verdadero y el dron no está ya armado ni volando, y así. Al pulsar los botones en la demo se sigue el ciclo real de un dron, cosa que permite ver visualmente su correcto funcionamiento.

El segundo widget es GestureDetector combinado con AnimatedContainer, patrón que usar los botones del HUD de la pantalla Tello, en los botones de flip y en los chips de modo. La sección muestra tres botones con estados enabled/disableed y active/inactive. AnimatedContainer solamente anima el color de fondo y el grosor del borde, haciendo que transicione suavemente entre estados.

El tercer widget es el botón de icono compacto (GestureDetector sobre un Container), usado para los botones de captura de foto, inicio/parada de grabación y cambio de cámara y mapa. A diferencia del widget anterior no usa AnimatedContainer, ya que el cambio de estado en estos botones es instantáneo. La sección muestra el botón de grabación con un toggle, donde al pulsarlo, cambia el icono y el color.

El cuarto widget documenta el patrón de presión continua, mediante GestureDetector con onLongPressStart, onLongPressEnd, onTapDown, onTap y onTapCancle, empleado en los botones de ajuste de la altitud en modo IMU. La combinación de los cinco callback cubre todos los posibles casos. 

El quinto widget es el Slider de zoom de la cámara, con el tema personalizado que usa SliderThemeData. La sección muestra el slider funcional con la etiqueta actualizada según el zoom.

El sexto widget es el DropdownButton del selector de modo de detección YOLO. La nota documenta los parámetros necesarios para que el desplegable encaje visualmente en la superposición sobre el vídeo.

El séptimo widget es el chip de modo, ModeChip y ConnectionChip, con AnimatedContainer. La sección muestra los selectores de modo de control (Classic, Voice, IMU) y de modo de dron (ArduPilot, SITL, Tello) completamente funcionales, con transición entre estados.

El octavo widget es el ModalBottomSheet combinado con DraggableScrollableSheet, usado en la pantalla de ayuda y en el selector de modos de la pantalla Tello. La sección incluye un botón que abre el sheet, mostrando el comportamiento de arrastre con scroll.

El noveno widget es el AlertDialog de confirmación. En este caso para la confirmación de eliminación.

El décimo widget es el ReordenableListView con handles de arrastre manuales mediante ReordenableDragStartListener, usado en la FlightPlanScreen en la lista de waypoints. La sección muestra una lista de cuatro waypoints de ejemplo completamente reordenables, con borrado individual. La nota documenta el detalle crítico del reordenamiento, que es necesaria por la forma en que Flutter calcula los índices tras eliminar el elemento de su posición original.

\subsection{Integración en EZDrone}
Esta demostración no se integra en EZDrone como funcionalidad, sino que actúa como agrupación de los diferentes elementos de diseño de la aplicación. Todos los widgets documentados se usan directamente en EZDrone con los mismos parámetros, colores y comportamientos descritos.

\subsection{Demostración}

\section{Joystick viruales}

\subsection{Caso de uso}

El mando RC de doble palanca, es el sistema de control de dron más usado entre los pilotos. Control cuatro canales de forma independiente: throttle y yaw en la palanca izquierda, pitch y roll en la derecha. Replicar esta disposición sobre una pantalla táctil convierte un móvil en un mando RC sin ningún tipo de hardware adicional.

El objetivo de esta demostración es documentar cómo se comportan los joysticks virtuales, rango de salida, comportamiento al soltarlo o mapeo a los cuatro ejes de vuelo, antes de integrarlo en una aplicación completa. Entender cómo funcionan esos valores es imprescindible para poder hacer una correcta conversión a señales PWM. 

\subsection{Implementación}
La demostración utiliza el paquete flutter-joystick (versión 0.0.4), que proporciona un widget “Joystick” que rastrea la posición del dedo relativamente al centro, y expone los valores continuos X e Y mediante un callback “listener”. La interfaz presenta dos joysticks configurados con “mode: JoystickMode.all”, que permite movimiento libre en los 360º, siguiendo la disposición de un control RC estándar.

Cada eje produce un valor en el rango [-1.0, 1.0], donde 0.0 representa la posición de reposo central. Una característica esencial del widget, es que cuando sueltas el dedo, el stick vuelve automáticamente al centro y el listener emite (0.0, 0.0). Esta característica es fundamental para el contexto de control de drones, sobre todo en modos como son el ALT-HOLD o el LOITER, donde un valor cero significa detener el movimiento.

Los valores de los cuatro ejes se muestran en pantalla en tiempo real, para que el desarrollador pueda verificar el rango de salida y la dirección antes de implementarlos en una aplicación. La conversión de estos valores normalizados a una señal PWM se realiza mediante la siguiente fórmula:

PWM = 1500 + [valor] * 500

Esta expresión mapea un valor de -1.0 a 1000 µs (deflexión máxima izquierda o abajo), 0.0 a 1500 µs (posición neutra) y +1.0 a 2000 µs (deflexión máxima derecha o arriba), cubriendo exactamente el rango PWM estándar de señales RC que espera el autopiloto.

\subsection{Integración en EZDrone}
En EZDrone, el mismo widget Joystick y la misma fórmula de conversión se utilizan en la pantalla de vuelo clásica (FlightScreen). La diferencia respecto a la demostración es la capa de transmisión. En lugar de mostrar los valores en la pantalla, los cuatro ejes se codifican en formato JSON ( {lx, ly, rx, ry} ) y se envían a la estación de tierra Python a través de un DataChannel WebRTC cada 50 ms (20Hz). La estación tierra recibe este paquete en el callback “on-joystick”, convierte los valores a PWM usando la misma fórmula y los envía al autopiloto mediante “dron.send-rc(roll, pitch, throttle, yaw)”.

Sobre la demostración base, EZDrone añade dos extras, el primero es una zona muerta de  ±0.1, donde los valores del joystick que estén en esa zona, se mapean automáticamente a 1500 µs, evitando que pequeñas imprecisiones generen un movimiento indeseado en el dron. El segundo es una curva de exposición aplicada al eje de yaw, añadida para suavizar la respuesta en pequeños movimientos, ya que el giro en el eje yaw es muy sensible a pequeñas variaciones.

\subsection{Demostración}

\section{Control por voz}
\subsection{Caso de uso}
El control por voz representa un paradigma de control de drones que libera las manos del operador durante el vuelo. En escenarios de inspección o monitorización, el operador puede necesitar tomar notas, manipular equipos o simplemente mantener la vista fija en el dron sin apartar la atención de la pantalla. Un sistema de voz bien diseñado permite emitir comandos verbalmente con el móvil a un lado.

El reto principal es el diseño de una interfaz de voz para el control de drones es su robustez, un sistema que reaccione a cualquier palabra detectada, generará comandos falsos de forma constante en cualquier ambiente ruidoso. La solución adoptada en esta demostración es una arquitectura de wake word, donde el sistema escucha pasivamente todo el rato, pero solo procesa comandos tras detectar la palabra de activación “dron”. Esto establece un umbral de activación sin requerir que el operador pulse ningún botón.

\subsection{Implementación}
La demostración utiliza el paquete speech-to-text (versión 6.6.2), que actúa como puente hacia el motor de reconocimiento de voz nativo del dispositivo: Google Speech en Android y el motor en Apple en iOS y macOS. Todo el procesamiento de reconocimiento es local, sin llamadas APIs externas, cosa que hace que se elimine la latencia de la red, y permite el funcionamiento sin conexión a internet una vez cargado el modelo.

El sistema de reconocimiento se modela como una máquina de estados de cuatro estados, implementada mediante el enum ListenState:

\begin{itemize}
    \item \textbf{idle} El sistema está inactivo, y el micrófono  apagado.
    \item \textbf{waitingWakeWord}  El micrófono está activo, esperando a que se diga la palabra de activación “dron”. 
    \item \textbf{waitingCommand}  La wake word ha sido detectada y se está esperando algún comando de vuelo.
    \item \textbf{commandLocked} El comando se ha ejecutado con éxito y el sistema permanece activo indefinidamente para recibir nuevos comandos.
\end{itemize}

Esta máquina de estados controla la lógica de reconocimiento y el estado visual de la interfaz, de forma que el operador siempre tiene una referencia visual de en qué estado se encuentra el sistema en todo momento.

El bucle de escucha se mantiene activo de forma continua mientras el sistema no está en estado idle. El motor de reconocimiento se reinicia automáticamente en los callback onStatus y onError para evitar que un silencio o un error puntual interrumpa la escucha. El parámetro pauseFor se establece en 4 segundos, para evitar que el motor corte comandos compuestos como podrían ser “mover adelante” o “volver al despegue”. El idioma de reconocimiento se puede cambiar, pero en esta demostración se fija en español “es-ES”.

Para gestionar el caso en que el operador olvida desactivar el sistema, se implementa un temporizador de apagado automático de 10 segundos. Este temporizador se resetea cada vez que se detecta cualquier voz y se cancela permanentemente cuando se ejecuta un comando con éxito, haciendo que el sistema se quede activo mientras el operador lo está usando.

Los comandos reconocidos se implementan en el método -identifyCommand, y cubren el ciclo de vuelo completo:

"armar"             -	ARM

"despegar"       	-	TAKEOFF

"subir"		        -	Ajuste de altitud positivo

"bajar" 	        -	Ajuste de altitud negativo

"mover adelante" 	-	Navegación hacia adelante

"mover atrás"  	    -	Navegación hacia atrás

"mover derecha" 	-	Navegación lateral derecha

"mover izquierda"	-   Navegación lateral izquierda

"aterrizar"  	    -   LAND

"volver a despegue" -   RTL

\subsection{Integración en EZDrone}
La integración del control por voz en EZDrone introduce una diferencia arquitectónica respecto a la demostración, ya que Flutter Web no puede utilizar el paquete speech-to-text (porque este depende de APIs nativas de la plataforma que no están disponibles en navegador). La solución que se adopta es la Web Speech API, soportada por todos los navegadores, que expone un puente  de interoperabilidad de JavaScript a Dart.

Esta capa se implementa en la clase WebSpeech, ubicada en el frontend Flutter. La clase define tipos de extensiones Dart, sobre los objetos JS correspondientes, lo que permite manipularlos desde Dart. 

En EZDrone la máquina de estados de cuatro fases de la demostración se sustituye por un modelo push-to-talk, más directo y adecuado para el uso de control de un dron. El operador mantiene pulsado el botón de micrófono, dicta el comando y lo suelta. Al soltar se llama al comando -speech.stopListening( ), y el texto reconocido se procesa. Este modelo elimina la necesidad de usar una wake word, ya que solo se dictan comandos cuando el operador lo decide.
\subsection{Demostración}

\section{Control inercial mediante IMU}

\subsection{Caso de uso}
La mayoría de los teléfonos móviles tienen una unidad de medida inercial (IMU) compuesta por acelerómetro, giroscopio y magnetómetro, cuya salida describe la orientación física del dispositivo en un espacio tridimensional. Mapear esta orientación a los ejes de control del dron crea una forma de controlarlo intuitiva, ya que inclinar el teléfono hacia adelante hace que el dron avance, inclinarlo hacia lateralmente lo desplaza hacia un lado y así. El operador sujeta el dispositivo con una empuñadura natural, y el control se genera mediante la postura y no mediante movimientos precisos de dedo.
Este paradigma es especialmente interesante para operadores menos familiarizados con la disposición de joysticks RC.

El reto técnico principal es que Flutter Web no puede acceder a los sensores del dispositivo directamente. El único camino disponible para ello es usar la API DeviceOrientationEvent del navegador, que requiere contexto HTTPS y en iOS, una solicitud explícita de permiso vinculada a un gesto del usuario.

\subsection{Implementación}
La solución se centra en torno a un puente JavaScript definido en el fichero index.html. Este puente registra un listener sobre el evento deviceorientation que va almacenando continuamente los tres ángulos de Euler en variables globales, y expone cinco funciones al entorno: “getAlpha( )” yaw, rango [0º, 360º], “getBeta( )”: pitch, rango [-180º, 180º],  “getGamma( )” roll, rango [-90º, 90º],  requestMotionPermission( ), para la compatibilidad con iOS, y finalmente stopOrientation( ), que elimina el listener.

En el lado de Dart, estas funciones se declaran mediante la anotación @JS, que permite llamar a funciones externas:

@JS('getBeta')  external double getBeta();
@JS('getGamma') external double getGamma();
@JS('getAlpha') external double getAlpha();

Un timer de 50 ms, llama continuamente a estas funciones, y se actualizan las variables internas de pitch, roll y yaw. Estos valores se muestran en tiempo real en tres recuadros. El botón “Activar sensores” aprovecha para combinar el inicio del timer con la petición obligatoria que requiere iOS.

Un aspecto práctico documentado en la demostración es que DeviceOrientationEvent requiere HTTPS, mientras que el servidor de Flutter sirve sobre HTTP. Para probar la demostración en un dispositivo físico, es necesario exponer el servidor local a través de un túnel HTTPS, la demostración incluye instrucciones paso a paso para hacerlo mediante ngrok.

\subsection{Integración en EZDrone}
En EZDrone, el mismo puente JavaScript, y el mismo timer se utilizan directamente en ImuFlightScreen. La integración añade dos elementos respecto a la demostración base.

El primero es el tratamiento de zona muerta.  Los ángulos brutos de los sensores no se pueden enviar directamente, ya que incluso con el dispositivo quieto, los valores fluctúan ligeramente, lo que generaría movimientos continuos. La función normalForward( double beta) implementa una zona muerta (entre -5º y 15º) para el eje del pitch, dado que el dispositivo en posición de reposo natural suele estar en esa inclinación. 
La función normalLateral (double gamma), hace lo mismo para el movimiento lateral con una zona muerta de ±10º. Una vez fuera de la zona muerta, se aplica una rampa lineal de 40º donde se mapean los valores para poder transformarlos a velocidades.

El segundo es el bloqueo de orientación de la pantalla. Cuando el dispositivo se inclina para controlar el dron, el sistema operativo puede girar la interfaz al detectar que el dispositivo se ha girado de más, haciendo la pantalla de EZDrone inutilizable. Esto se soluciona llamando a sceen.orientation.lock( ) mediante el puente JS definido en html.index. Al salir de la pantalla, unlockOrientationEZ( ) libera este bloqueo.

\subsection{Demostración}

\section{Localización GPS en mapa}

\subsection{Caso de uso}
Conocer la posición del operador en tiempo real es una funcionalidad con mucho valor operativo, permite ver la distancia real entre el operador y el dron, mostrar ambas posiciones en el mismo mapa, y en un escenario futuro, implementar modos de seguimiento o establecer zonas de exclusión relativas a la posición del operador. 

Esta demostración documenta cómo obtener la posición GPS del dispositivo y visualizarla en tiempo real sobre un mapa de código abierto.

\subsection{Implementación}
La demostración utiliza dos paquetes: geolocator (versión 13.0.2) para la obtención de la posición GPS, y flutter-map (versión 7.0.2) para el renderizado del mapa, con latlong2 (versión 0.9.1) para el manejo de coordenadas. La elección de flutter-map sobre alternativas es intencional, ya que no requiere claves API y es completamente abierta. Utiliza tiles de OpenStreetMap, lo que la hace adecuada para el contexto del proyecto, y elimina dependencias de servicios de pago.

La posición se obtiene mediante Geolocator.getPositionStream( ), que abre un stream que emite actualizaciones cada vez que el dispositivo se desplaza al menos 1 metro (configurable). Cada posición nueva se acepta siempre si es la primera (para evitar un bucle de espera) o si la precisión es inferior a 50 metros. Las lecturas con mayor precisión se descartan para evitar saltos bruscos en el mapa. 
La precisión GPS actual se muestra en un badge superpuesto sobre el mapa con una codificación de colores, verde por debajo de 5 m, naranja por debajo de 20 m y rojo por encima.

\subsection{Integración en EZDrone}
En EZDrone la misma lógica de stream GPS se usa en los métodos startUserLocation( ) y stopUserLocation( ) del DronProvider. La posición del operador se almacena en userPosition, y se inicia automáticamente al conectar al dron, deteniéndose al desconectar. La posición se visualiza dentro del mapa como un icono de teléfono móvil azul, permitiendo al operador percibir de forma visual la separación entre él y el dron.

\subsection{Demostración}

\section{Streaming de vídeo WebRTC con detección YOLO}

\subsection{Caso de uso}
La visión en primera persona (FPV) es una capacidad fundamental en cualquier estación de control de drones para misiones de inspección, fotografía aérea o navegación en entornos complejos. Esta demostración documenta cómo establecer un stream de vídeo en tiempo real desde una cámara conectada al ordenador. Se utiliza WebRTC como protocolo de transporte, y se integra detección de objetos en tiempo real mediante YOLO sobre el stream antes de su transmisión.

\subsection{Implementación en el emisor Python}
El emisor es un script Python, que utiliza las bibliotecas aiortc y websockets. Define una clase CustomVideoStreamTrack que hereda de VideoStreamTracl de aiortc y sobreescribe el método asíncrono recv( ).

En cada llamada a recv( ) (que es en cada frame que se tiene que transmitir), el track captura un frame de la cámara mediante OpenCV, ejecuta interferencia YOLO cada 25 frames, dibuja los bounding boxes sobre el frame según el modo de detección activo, añade un timestamp y convierte el frame al formato rgb24 requerido por aiortc.

El servidor WebSocket permanece abierto durante toda la sesión, actuando tanto de servidor de señalización como canal de control en tiempo real para los comandos que Flutter envía durante el streaming (cambios de modo de detección: all/person/none, solicitud de captura de frame e inicio/parada de grabación).

\subsection{Implementación en el cliente Flutter}
El cliente Flutter utiliza la librería flutter-webrtc. Al conectar, crea una RTCPeerConnection con el servidor STUN de Google, registra un callback onTrack que asigna el stream de vídeo recibiendo al renderer RTCVideoRenderer, y abre el canal WebSocket hacia el emisor Python.

La captura de pantalla y la grabación de vídeo se realizan en el cliente mediante interoperabilidad JS, sin enviar ningún dato al servidor Python. Las funciones captureDroneFrame( ) y startDroneRecording( ) / stopDroneRecording( ), definidas en el index.html, acceden al stream y lo extraen mediante captureStream( ) y utilizan la API MediaRecorder del navegador para grabar en formato WebM. Dart llama a estas funciones mediante @JS y la biblioteca dart:js-interop.

\subsection{Integración en EZDrone}
En EZDrone, el stream WebRTC se recibe en WebRTCLogic y el MediaStream resultante se expone al DronProvider mediante el callback onRemoteStream. La pantalla de vuelo clásica puede alternar entre el mapa satelital y el stream de vídeo mediante el botón de intercambio, renderizando el stream mediante el widget RTCVideoView. Las mismas funciones JS de captura y grabación se invocan desde el DronProvider mediante jsCapture( ) y jsStartRecording( ) / jsStopRecording( ), exponiendo estas funciones como botones en la interfaz de vuelo. El modo de detección YOLO se controla desde la misma pantalla mediante un widget dropdown que deja elegir entre personas, todo y nada. Entonces, mediante el topic MQTT mobileFlutter/groundStation/detectionMode, Flutter envía la elección a la estación tierra, donde se produce el cambio de detection-mode.

\subsection{Demostración}

\section{Planificación de misiones autónomas}

\subsection{Caso de uso}
Los paradigmas de control descritos en los apartados anteriores son todos reactivos, es decir, el operador emite comandos en tiempo real y el dron responde a ellos de forma inmediata. Existen una clase de operaciones donde lo que se busca es definir la misión con antelación y dejar su ejecución al autopiloto, sin requerir de la atención e intervención continua del operador (como podrían ser inspecciones de infraestructuras, cartografías aéreas, vigilancia de perímetros, riego de cultivos entre muchos otros).

ArduPilot soporta este modo de vuelo autónomo, permite que se le pasen una lista de waypoints con coordenadas GPS y altitudes, los almacena en la memoria y los ejecuta en modo AUTO.

Esta demostración abarca el problema de diseñar una interfaz de planificación de misiones accesible desde cualquier navegador. A parte de demostrar cómo construir una lista de waypoints de forma visual sobre un mapa, asignar acciones a cada uno de los waypoints, guardar los planes entre sesiones, exportarlos a un formato compatible con herramientas como MissionPlanner y enviarlos al dron para su ejecución.

\subsection{Modelos de datos}
La demostración describe el modelo de datos de la misión, que en EZDrone se implementa en tres clases distintas del Provider.

FlightWaypoint representa un punto de paso individual. Sus campos son lat y lon (coordenadas en grados), altM (altitud sobre el suelo en metros) y action, un objeto WaypointActoin. Este objeto contiene un tipo de acción a ejecutar al llegar al waypoint junto a su duración si es el caso. Los tipos de acciones disponibles son: none (volar sobre el punto sin detenerse), hover (mantener posición durante los segundos definidos), takePhoto (disparar la cámara), recordVideo (grabar durante el tiempo definido), rtl (retorno al punto de lanzamiento) y land (aterrizaje en el waypoint).

FlightPlan agrupa una lista de FlightWaypoints con un identificador único, un nombre editable y una marca de tiempo de creación. Esta clase también contiene un método para calcular las distancias entre waypoints consecutivos, totalDistanceM, usando la fórmula de Haversine: 

const R = 6371000.0;
final s = sin(dLat/2)*sin(dLat/2) + cos(lat1)*cos(lat2)*sin(dLon/2)*sin(dLon/2);
d += R * 2 * atan2(sqrt(s), sqrt(1 - s));

La elección de Haversine es más que correcta para esta escala, para misiones de varios kilómetros, la curvatura terrestre introduce errores no despreciables si se usa geometría plana.

\subsection{Implementación}
La demostración presenta una interfaz de planificación de misiones completamente funcional sin necesidad de dron ni simulador. 

La pantalla de planificación presenta una interfaz separada en dos zonas. La zona superior es un mapa satelital de flutter-map. Un tap sobre cualquier punto del mapa llama al callback de onTap, que añade un nuevo FlightWaypoint en el punto donde se ha pulsado. Los waypoints se ven como marcadores numerados, el primero en verde (indicando el inicio) y el último en rojo (indicando el final). Una línea lila los une en orden, y un badge superpuesto muestra el número de waypoints y la distancia total calculada.

La zona inferior es una lista de Waypoints, implementada usando ReordenableListView, que permite reordenar los puntos mediante arrastrar y soltar. Cada elemento de la lista muestra las coordenadas con cinco decimales, la altitud asignada, el chip de acción si existe y un botón para borrarlo. Al pulsar sobre un elemento se abre un AlertDialog de edición, donde el operador puede modificar la altitud y seleccionar el tipo de acción y su duración.

Los planes se guardan en el almacenamiento local del navegador, usando web.window.localStorage. La serialización y deserialización se realiza mediante los métodos toJson( ) y fromJson( ). 

La demostración incluye exportación a formato QGC, el estándar de misiones en el ecosistema de ArduPilot. Se genera un fichero de texto plano con todos los waypoints y su comando MAVLink correspondiente para sus acciones. La descarga se hace en el propio navegador usando la API Blob.

\subsection{Integración en EZDrone}
En EZDrone, la demostración se extiende con dos acciones adicionales, el envío del plan al dron y su ejecución. El botón UPLOAD serializa el plan como JSON y lo publica en el topic MQTT /uploadMission, la estación tierra lo almacena y confirma su recepción. Al pulsar START, la estación ejecuta el plan mediante dron.goto( ). Durante la ejecución, el mapa muestra el progreso con el trail y el waypoint activo resaltado comparado con los otros.

\subsection{Demostración}

\section{Registro y análisis de vuelos}

\subsection{Caso de uso}
El registro de los vuelos realizados tiene valor tanto operativo como de revisión. Conocer la trayectoria recorrida, la evolución de la carga de la batería, la velocidad pico alcanzada o el tiempo transcurrido desde el despegue, permite al operador evaluar el rendimiento del vuelo, detectar comportamientos anómalos y documentar las operaciones. Esta demostración aborda el diseño de una pantalla de historial de vuelos que combina visualización de trayectorias, reproducción temporal de la telemetría y exportación de los datos.

El objetivo de esta demostración es en particular doble, por un lado, documentar el modelo de datos y la lógica de visualización de forma aislada sin depender de un dron real, y por otro lado, servir como prueba antes de implementarlo en EZDrone con datos reales.

\subsection{Datos de ejemplo}
Dado que la demostración se ejecuta sin ningún tipo de dron o SITL, los datos han sido generados manualmente en el fichero sample-flight-data.dart. Este fichero contiene tres sesiones de ejemplo con coordenadas reales sobre el campus. La primera se trata de un vuelo cuadrado a 20 m en modo clásico (un tipo de modo de control de EZDrone) con SITL, de 4 minutos y 32 segundos, completado correctamente, la segunda es un vuelo triangular a 30 m en modo IMU (modo de control de EZDrone), de 6 minutos y 10 segundos, también completado. Finalmente la tercera se trata de un vuelo parcial a 15m en modo voz (modo de control de EZDrone), de 1 minuto y 48 segundos, interrumpido por desconexión. Esta combinación de sesiones con distintos modos de control, flags y estados de finalización, permite verificar diferentes opciones posibles de la interfaz.

Aparte, se ha implementado la función GenLog( ), capaz de generar la telemetría de cada sesión de forma coherente con la trayectoria. La altitud sube linealmente durante el primer 15/100 del vuelo, se mantiene con pequeñas fluctuaciones aleatorias, y en el último 15/100 baja. La batería decae de forma lineal, la velocidad es baja en los tramos de despegue y aterrizaje y se estabiliza entre 3 y 5.5 m/s en crucero. El rumbo en cada punto se calcula mediante trigonometría sobre el vector entre el punto actual y el siguiente, reproduciendo la orientación real del dron en cada segmento.

\subsection{Implementación}
La pantalla principal de la demostración, FlightLogScreen, presenta la lista de sesione en un ListView. Cada sesión se muestra en una tarjeta que incluye la fecha y hora de inicio, un chip de color con el modo de control usado, un chip naranja SITL cuando el vuelo ha sido simulado, un indicador de estado completado o interrumpido y cinco métricas calculadas: duración, distancia total, altitud máxima, desnivel y velocidad máxima.

Al pulsar sobre las tarjetas, se abre un bottom sheet, implementado mediante un DraggableScrollableSheet. La zona superior contiene un mapa con la trayectoria completa dibujada. Tres marcadores que identifican el punto de despegue (verde), el de aterrizaje / interrupción (rojo / naranja) y la posición actual del dron en cada momento.

El slider temporal en el elemento central de la bottom sheet. Su rango va de 0 al número de puntos del trail, y al arrastrarlo se actualizan simultáneamente el segmento visible del trail, la posición del dron y el texto de telemetría. 

Debajo del slider, se presenta una cuadrícula de ocho tarjetas que aparte de incluir las métricas de la tarjeta de resumen, incluyen la batería mínima registrada y el número total de snapshots del log. El botón DOWNLOAD CSV construye el fichero en formato CSV, y lo descarga en el navegador, pero en esta demostración se muestra el contenido en un diálogo si la descarga no está disponible.

\subsection{Integración en EZDrone}
En EZDrone, se usan directamente los datos obtenidos en el vuelo real o simulado, mediante el Provider. Al aterrizar, se cierra la sesión y se inserta en el flightHistory. La lista completa se guarda en el almacenamiento local. Y seguidamente se muestran en la pantalla de FlightLogScreen.

\subsection{Demostración}



\chapter{Evolución de la arquitectura}
La arquitectura presentada en el Capítulo 3 cumple su objetivo principal, demuestra que es posible controlar un dron ArduPilot desde un navegador usando Flutter. El problema aún así, son sus limitaciones estructurales. La arquitectura no llega a los requerimientos de una estación de control operativo. Este capítulo explica la evolución desde ese diseño inicial hasta la arquitectura híbrida definitiva adoptada en EZDrone, documentando los problemas técnicos que han motivado a este cambio y la implementación de sus soluciones.

\section{Limitaciones de la arquitectura base}
La arquitectura de tres capas del Capítulo 3 (Flutter Web, backend Dart y estación tierra Python) presenta tres problemas o limitaciones principales.

La primera es la latencia introducida por el polling HTTP. El frontend consulta constantemente el endpoint /status (cada 500 ms), para así obtener la telemetría y el estado del dron. Esto implica que cualquier evento puede tardar ese medio segundo en reflejarse en la interfaz. Para un sistema de monitorización a tiempo real, medio segundo puede llegar a ser aceptable, pero a la hora de controlar el movimiento del dron empieza a ser un problema serio, ya que el operador necesita una retroalimentación inmediata.

EXLICAR PROBLEMA PETICIONES HHTP

La tercera limitación, y la más estructural, es la incapacidad de soportar canales de alta frecuencia. El control por joystick, debe enviarse a 20 Hz (cada 50 ms) para poder proporcionar una respuesta relativamente fluida. Implementar esto mediante peticiones HTTP POST, generaría 1200 peticiones por minuto solamente para el joystick, cosa que saturaría tanto el backend como el broker MQTT. No solo eso sino que toda la latencia acumulada haría imposible el control del dron. El streaming de vídeo es también es directamente inviable sobre HTTP.

A estas limitaciones técnicas se añade una de diseño, el backend Dart actúa como proxy entre Flutter y el broker MQTT, añadiendo así un salto de red y un punto extra de posible fallo, y todo esto sin añadir ninguna cosa que no se pudiera implementar directamente en el propio frontend.

\section{Primera evolución: eliminación del backend Dart}
El primer paso de esta evolución es eliminar el backend Dart como intermediario y conectar Flutter directamente al broker MQTT. Esto es posible gracias al soporte de WebSocket en HiveMQ, aunque los navegadores no permiten conexiones TCP directamente, sí que permiten conexiones WebSocket.

La implementación se asienta en una clase llamada MqttLogic, directamente ubicada en el frontend. A diferencia de la arquitectura anterior que usaba MqttServerClient en el backend, esta clase inicia MqttBrowserClient con la url del broker (wss://broker.hivemp.com/mqtt) y el puerto 8884. La variante segura del WebSocket, WSS es necesaria porque la página HTTPS no permite abrir conexiones WS no seguras y el navegador directamente las bloquea.

Esta clase contiene tres operaciones principales, connect( ), que establece la conexión y registra el listener de mensajes entrantes, subscribe(topic), para suscribirse a un topic y publish(topic, payload), para publicar. Los mensajes entrantes se entregan al DronProvider, desacoplando así la capa de transporte de la lógica interna de la aplicación.

Con este cambio, el flujo de un comando de ARM (por ejemplo), pasa de cuatro saltos: Flutter - HTTP - Dart - MQTT - Python, a solamente dos: Flutter - MQTT - PYTHON. El estado de la sesión deja de depender de un proceso Dart local, y la aplicación puede desplegarse como un conjunto de archivos estáticos sin necesidad de mantener un servidor intermediario en marcha.

La gestión del estado sigue siendo el mismo modelo de estados del Capítulo 3, pero ahora el Provider es quien reacciona directamente a los topics MQTT suscritos, en lugar de tener que hacer un polling constante. Cuando la estación tierra publica, el método handleMessage del provider reacciona.

\section{Segunda evolución: Incorporación de WebRTC}
La eliminación del backend como intermediario resuelve la latencia del polling, pero no aborda nada de los canales de alta frecuencia. Para el joystick, la telemetría continua y el vídeo, se añade WebRTC, donde su arquitectura peer-to-peer, elimina el broker como figura intermediaria y proporciona la latencia y el ancho de banda necesario para estos canales.

La arquitectura de comunicación definitiva asigna a cada protocolo los canales para los que es más adecuado, basándose en las características discutidas en el Capítulo 2.

MQTT sobre WebSocket se mantiene para todos los comandos discretos (connect, arm, takeoff, land, rtl, disconnect, setMode, setCamera, detectionMode, uploadMission, startMission) y para las confirmaciones de cambio de estado. Estos mensajes son poco frecuentes, y su pérdida sería crítica.

WebRTC se usa para los tres canales mencionados anteriormente. Un DataChannel para el joystick, donde se envían los valores de los cuatro ejes desde Flutter a la estación tierra. La baja latencia de WebRTC es imprescindible para esto, ya que un retardo entre el movimiento del joystick y la respuesta del dron podría provocar problemas de control.

Otro DataChannel para la telemetría, que envía la telemetría desde la estación tierra hacia Flutter. 

El tercero es un video track que transmite el stream de la cámara del dron desde Python a Flutter. El vídeo es el caso más evidente de uso de WebRTC. No tiene alternativa razonable ya que requiere una actualización demasiado rápida para llevarse sobre MQTT o HTTP.

\section{Servidor de señalización}
La conexión WebRTC requiere un servidor de señalización para el intercambio inicial de SDP y candidatos ICE entre los dos peers. Este servidor se implementa en el backend, que queda desplegado en el propio servidor de la eetac: dronseetac.upc.edu.

El servidor es una aplicación Dart pura que levanta un servidor HTTPS sobre el puerto 8105, cargando los certificados TLS de Let’s Encrypt. Atiende dos tipos de peticiones, las que llegan a la ruta /ws se tratan como conexiones WebSocket de señalización. El resto sirven como ficheros estáticos a la carpeta /web, que contiene el build de Flutter Web compilado. Esta arquitectura permite que el mismo proceso sirva tanto la aplicación al navegador como el canal de señalización, sin necesidad de un servidor web separado.

El protocolo de señalización implementado tiene tres fases. En la primera, el cliente envía su identificador de peer “browser” para Flutter y “Python” para la estación tierra. El servidor registra el par y responde inmediatamente con la configuración ICE, que incluye un servidor STUN público y el servidor TURN privado (alojado en dronseetac.upc.edu:3478). 

En la segunda fase, los mensajes posteriores (SDP offer/answer y candidatos ICE) se reenvían al peer indicado. En la última, cuando un peer se desconecta, se elimina su entrada.

Explicar que es un servidor TURN????

\subsection{Certificación TLS}
Para que el servidor Dart pueda atender conexiones HTTPS y WSS, es necesario que disponga de un certificado TLS válido y reconocido por los navegadores. En este proyecto, el certificado se obtiene mediante Let's Encrypt, una Autoridad Certificadora gratuita y de confianza.

El proceso de emisión del certificado utiliza la herramienta certbot, que se ejecuta directamente en el servidor (en este caso en dronseetac.upc.edu). El flujo de certificación sigue los pasos siguientes:

certbot solicita un certificado a Let’s Encrypt para el dominio dronseetac.upc.edu. Let’s Encrypt emite un challenge de tipo HTTP-01, donde exige que el servidor sirva un fichero de verificación en la ruta /.well-known/acme-challenge/<token>. Certbot despliega ese fichero y notifica a Let’s Encrypt para que lo verifique. Una vez superada la verificación, Let’s Encrypt firma y emite el certificado.

El resultado del proceso son dos ficheros que certbot almacena en el sistema, fullchain.pem, la cadena completa de certificados del servidor y los certificados intermedios de la CA. Y privkey.pem, la clave privada del servidor.

Los certificados de Let’s Encrypt tienen una validez de 90 días. Para evitar interrupciones, certbot instala automáticamente una tarea que intenta la renovación cuando el certificado está a punto de caducar. Tras la renovación, los nuevos certificados reemplazan los anteriores y el servidor se reinicia para cargar los nuevos certificados. 

\section{Capa WebRTC en Flutter}
La lógica WebRTC del cliente Flutter, se encuentra dentro de la clase WebRTCLogic, en el frontend. Esta clase gestiona el ciclo de vida de la conexión, es decir, el establecimiento de la conexión, recepción de stream de vídeo, apertura de los DataChannels y desconexión.

Al llamar a connect(signalingUrl), la clase abre una conexión WebSocket hacia el servidor de señalización y envía su identificador (“browser”) como primer mensaje. A continuación, espera a la configuración ICE. Con la configuración ICE recibida, crea la RTCPeerConnection.

\section{Arquitectura híbrida definitiva}
La arquitectura resultante combina los dos protocolos de forma que cada uno opera exclusivamente en el dominio para el que es más adecuado. La arquitectura final tiene la siguiente estructura:

AÑDIR DIAGRAMA

Esta separación tiene muchas ventajas importantes. El broker MQTT solo recibe tráfico de baja frecuencia, joysticks y telemetría fluyen directamente entre los peers una vez se ha establecido la conexión WebRTC, sin tener que pasar por ningún servidor intermediario.

Una implicación operativa importante es el orden de inicio, donde la estación tierra debe iniciar la conexión MQTT al broker antes de poder recibir el comando connect, y debe de estar el canal WebRTC abierto antes de poder recibir datos o enviar telemetría. 




\chapter{EZDrone: La estación de tierra integrada}
Los capítulos anteriores han recorrido el camino desde la arquitectura original basada en HTTP hasta la arquitectura híbrida MQTT/WebRTC descrita en el capítulo 7. Este capítulo presenta finalmente EZDrone, el sistema resultante de integrar esa arquitectura de comunicaciones con una aplicación de control completa. Mientras que el Capítulo 7 se centra en cómo viajan los datos, este se centra más en qué ofrece la aplicación al operador, la gestión centralizada del estado, los tres modos de control, el soporte de tres tipos distintos de drones, la plataforma de misiones, la visualización en tiempo real y el registro de los vuelos. El capítulo cierra con una demostración integral del sistema en funcionamiento.

\section{Visión general de EZDrone}
EZDrone es una GCS que se ejecuta en el propio navegador. A diferencia de las soluciones de escritorio clásicas como Mision Planner o QGroundControl, no requiere instalación previa, la aplicación Flutter Web se sirve como un conjunto de ficheros estáticos a los que el operador accede desde cualquier dispositivo con cualquier navegador. Esta decisión de diseño es la que habilita los modos de control por voz e inercial, que se apoyan en sensores y APIs propias del navegador del dispositivo.

El sistema se asienta sobre la arquitectura híbrida del Capítulo 7, que reparte cada flujo de información por el protocolo más adecuado a sus características. Los comandos discretos y las confirmaciones de estado viajan por MQTT sobre WebSocket seguro, los flujos de alta frecuencia viajan por canales WebRTC peer-to-peer, que aportan la baja latencia y el ancho de banda que MQTT no puede ofrecer. 

Sobre esta base, EZDrone añade una capa funcional que es el objeto de este capítulo, un único punto de gestión del estado (el DronProvider), tres modos de control alternativos (clásico, por voz e inercial), soporte para drones (ArduPilot reales, el simulador SITL, y el DJI Tello), planificación y ejecución de misiones autónomas, visualización cartográfica y de vídeo en tiempo real, y un sistema de registro y galería multimedia persistente en el propio navegador.

\section{Gestión de estado: el DronProvider y la máquina de estados}
Toda la lógica de estado de la aplicación se concentra en una única clase, el DronProvider. Los widgets nunca almacenan el estado de vuelo, sino que se limitan a leerlo para redibujarse (\texttt{context.watch}) o a invocar acciones sobre él (\texttt{context.read}). Esta centralización es aposta, ya que garantiza que las distintas pantallas observen siempre una única fuente, y queden desacopladas de la lógica de transporte (MQTT y WebRTC).

El Provider mantiene tres tipos de estado. En primer lugar, un conjunto de enumeraciones que fijan la configuración de la sesión, \texttt{ControlMode} (clásico, voz, IMU), \texttt{DroneConnectionMode} (arduPilot, sitl, tello) y \texttt{DetectionMode} (all, person, none). En segundo lugar, una colección de variables booleanas que dirigen qué botones se habilitan y que muestra la interfaz (\texttt{isConnected, isArmed, isConfigValid, isMissionMode,} entre otras). Y en tercer lugar, las variables de telemetría (\texttt{currentLat, currentLon, currentAlt, currentBat, currentHeading...}), que se actualizan a partir del canal de telemetría WebRTC.

Estas variables no son independientes entre sí, sino que están ligadas por una máquina de estados que reproduce el ciclo de un dron real, (\texttt{deconectado - conectado - armado - volando - aterrizado}). Cada transición dispara una confirmación enviada por la estación de tierra, no por el propio comando. Es decir, al pulsar ARM, la aplicación no se considera armada hasta que recibe el mensaje \texttt{armed} de vuelta. Hasta entonces, permanece en un estado de espera. Esto evita que la interfazz se adelante al estado físico del dron.  

(tabla de botones y booleanos??????)

\section{Soporte multi-dron}
Una de las aportaciones de EZDrone respecto a la arquitectura original es el soporte de tres destinos de vuelo distintos, seleccionables desde la pantalla de configuración inicial sin cambios en el cliente. La elección se publica al conectar mediante el topic \texttt{setMode} y determina como la estación tierra establece el enlace.

Los modos \texttt{ardupilot} y \texttt{sitl} comparten toda la lógica de pilotaje, los mismos modos y funciones, basados en la librería dronLink, que implementa el protocolo MAVLink, y descompone las diferentes acciones (armar, despegar, ir a un punto, telemetría etc.) en diferentes módulos propios. La única diferencia entre ambos es el extremo de conexión, un autopilot Pixhawk en el caso de dron real, o el proceso del simulador SITL en el caso del simulado. Esta equivalencia en general, es la que permite desarrollar y validar el sistema sin necesidad de probarlo en el hardware real.

El modo \texttt{tello} recurre a una librería distinta, la de TelloLink, ya que DJI no usa MAVLink, sino su propio SDK sobre WI-FI. La pantalla del Tello en si también es una implementación separada (\texttt{TelloFlightScreen}), diseñada intentando replicar la propia pantalla de control de DJI en torno a los modos especiales que ofrece el tello y que se describen en el apartado 8.6.

\section{Modos de control (ArduPilot / SITL)}
EZDrone ofrece tres formas alternativas de pilotar el dron en los modos ArduPilot y SITL. Todas producen el mismo tipo de salida (valores normalizados de cuatro ejes enviados por el DataChannel del joystick), pero difieren radicalmente en la interfaz de entrada. El modo se elige en la pantalla de configuración y determina a qué pantalla de vuelo se accede.

\subsection{Modo clásico (FlightScreen)}
La pantalla clásica (FlightScreen) reproduce el control de un mando de rc mediante dos joysticks virtuales y un conjunto de instrumentos. En modo apaisado, la pantalla se divide en tres columnas, el panel izquierdo muestra instrumentos numéricos (altitud, velocidad, latitud y longitud), el panel central contiene el mapa satélite o el vídeo FPV con los dos joysticks superpuestos y el panel derecho añade la rosa de los vientos animada según la orientación, el propio heading y el estado del vuelo. En modo vertical, los instrumentos se compactan en una sola fila horizontal en la parte superior.

El stick izquierdo lleva el thrust y el yaw, mientras que el derecho el pitch y roll. Los ejes del stick izquierdo se atenúan por un factor de 0.35 para suavizar las entradas (ya que es un movimiento muy brusco). La barra de acción inferior contiene los botones ARM, TAKEOFF, LAND, RTL y DISCONNECT, cuya habilitación depende directamente de los booleanos del Provider descritos en el apartado 8.2.

Sobre el vídeo en directo, se superponen varios controles. Un dropdown para seleccionar el modo de detección YOLO (Todo/personas/nada), botones de captura de foto y grabación de vídeo, un selector de cámara (dron/portátil con ET) y una barra de zoom (de 1x a 5x) que se orienta verticalmente en modo portrait y horizontalmente en modo landscape. Cuando el vídeo se está grabando, aparece el badge de \texttt{REC} en la esquina durante toda la grabación
\input{memoria/chapters/chap9_sostenibilidad}
\input{memoria/chapters/chap10_conclusiones_y_futuro}


%%%  BIBLIOGRAFIA
%%%%%%%%%%%%%%%%%%%%%%%%%%%%%%%%%%%%%%%%%%%%%%%%%%%%%%%%%%%%%%%%%%%%%%%%%%

\begin{thebibliography}{99}

\bibitem{UAV applications}
G.N. Muchiri and S. Kimathi,
\textit{A Review of Applications and Potential Applications of UAV},
en \textit{Proceedings of the Sustainable Research \& Innovation (SRI) Conference},
Jomo Kenyatta University of Agriculture and Technology,
publicado el 04 abr. 2022.
Disponible en: \url{https://sri.jkuat.ac.ke/jkuatsri/index.php/sri/article/view/325}.

\bibitem{missionplanner}
M. Oborne,
\textit{Mission Planner Ground Control Station},
ArduPilot, open-source software, 2013.
Disponible en: \url{https://ardupilot.org/planner/}.
[Consultat: 04 juny 2026].

\bibitem{qgroundcontrol}
MAVLink / QGroundControl Dev Team,
\textit{QGroundControl: Cross-platform Ground Control Station for Drones},
open-source software.
Disponible en: \url{http://qgroundcontrol.com}.
[Consultat: 04 juny 2026].

\bibitem{flutter}
Google LLC,
\textit{Flutter -- Build apps for any screen},
2024.
Disponible en: \url{https://flutter.dev}

\bibitem{dronsEETAC}
dronsEETAC,
\textit{Drone Engineering Ecosystem (DEE)},
repositori GitHub, Universitat Politècnica de Catalunya -- EETAC, 2022.
Disponible en: \url{https://github.com/dronsEETAC/DroneEngineeringEcosystemDEE}.
[Consultat: 04 juny 2026].

\bibitem{pixhawk}
3DR / ArduPilot,
\textit{Pixhawk Autopilot Overview},
documentació tècnica, ArduPilot.
Disponible en: \url{https://ardupilot.org/copter/docs/common-pixhawk-overview.html}.
[Consultat: 04 juny 2026].

\bibitem{ardupilot}
ArduPilot Dev Team,
\textit{ArduPilot Open Source Autopilot},
open-source firmware.
Disponible en: \url{https://ardupilot.org}.
[Consultat: 04 juny 2026].


\bibitem{tello}
Ryze Technology,
\textit{Tello Drone -- Specifications},
Shenzhen Ryze Technology Co., 2018.
Disponible en: \url{https://www.ryzerobotics.com/es/tello}.
[Consultat: 04 juny 2026].

\bibitem{px4}
PX4 Dev Team,
\textit{PX4 Open Source Autopilot},
open-source firmware.
Disponible en: \url{https://px4.io}.
[Consultat: 04 juny 2026].

\bibitem{mavlink}
MAVLink Dev Team,
\textit{MAVLink Micro Air Vehicle Communication Protocol},
open-source protocol.
Disponible en: \url{https://mavlink.io}.
[Consultat: 04 juny 2026].

\bibitem{ardupilot_gcs}
ArduPilot Dev Team,
\textit{Choosing a Ground Station},
documentació oficial ArduPilot.
Disponible en: \url{https://ardupilot.org/copter/docs/common-choosing-a-ground-station.html}.
[Consultat: 04 juny 2026].


\bibitem{skydio}
Skydio,
\textit{Skydio Autonomy -- Autonomous Drone Technology},
Skydio, Inc.
Disponible en: \url{https://www.skydio.com/software}.
[Consultat: 04 juny 2026].


\bibitem{djifly}
DJI,
\textit{DJI Fly -- App},
DJI Technology Co., Ltd.
Disponible en: \url{https://www.dji.com/es/dji-fly}.
[Consultat: 04 juny 2026].


\bibitem{mqtt_dart}
Tim Carey-Smith et al.,
\textit{mqtt\_client -- A server and browser based MQTT client for Dart},
pub.dev, Dart package.
Disponible en: \url{https://pub.dev/packages/mqtt_client}.
[Consultat: 04 juny 2026].


\bibitem{sun2025}
S.-C. Sun and Y.-C. Tsai,
\textit{A Modular AIoT Framework for Low-Latency Real-Time Robotic
Teleoperation in Smart Cities},
arXiv preprint arXiv:2510.11421, oct. 2025.
DOI: \url{https://doi.org/10.48550/arXiv.2510.11421}.

\bibitem{hexsoon}
dronsEETAC,
\textit{Guia Hexsoon EDU-450: Muntatge, configuració i proves de vol},
repositori GitHub, Universitat Politècnica de Catalunya -- EETAC.
Disponible en: \url{https://github.com/dronsEETAC/GuiaHexsoonEDU450}.
[Consultat: 04 juny 2026].

\bibitem{telloEETAC}
dronsEETAC,
\textit{Tello Engineering Ecosystem: A software tool for controlling
a Tello drone using different devices and applications},
repositori GitHub, Universitat Politècnica de Catalunya -- EETAC, 2023.
Disponible en: \url{https://github.com/dronsEETAC/TelloEngineeringEcosystem}.
[Consultat: 04 juny 2026].

\bibitem{dart}
Google,
\textit{Dart Programming Language},
open-source, Google LLC, 2011.
Disponible en: \url{https://dart.dev}.
[Consultat: 04 juny 2026].


\bibitem{dart2js}
Google,
\textit{dart2js: Dart-to-JavaScript compiler},
documentació oficial Dart.
Disponible en: \url{https://dart.dev/tools/dart2js}.
[Consultat: 04 juny 2026].

\bibitem{flutter_arch}
Google,
\textit{Flutter Architectural Overview},
documentació oficial Flutter.
Disponible en: \url{https://docs.flutter.dev/resources/architectural-overview}.
[Consultat: 04 juny 2026].

\bibitem{provider}
R. Rousselet,
\textit{provider: A wrapper around InheritedWidget to make them
easier to use and more reusable},
pub.dev, versió 6.1.5+1, Dart package, 2024.
Disponible en: \url{https://pub.dev/packages/provider}.
[Consultat: 04 juny 2026].

\bibitem{hivemq}
HiveMQ GmbH, ``HiveMQ: The Real-Time Industrial Data Platform,''
[Online]. Available: \url{https://www.hivemq.com}.
[Accessed: Jun. 5, 2026].

\bibitem{mosquitto}
Eclipse Foundation, ``Eclipse Mosquitto: An Open Source MQTT Broker,''
[Online]. Available: \url{https://mosquitto.org}.
[Accessed: Jun. 5, 2026].

\bibitem{ardupilot_sitl}
ArduPilot Dev Team, ``Mission Planner Simulation (SITL),''
[Online]. Available: \url{https://ardupilot.org/planner/docs/mission-planner-simulation.html}.
[Accessed: Jun. 4, 2026].

\bibitem{api}
IBM, ``What is an API (Application Programming Interface)?''
IBM Think, Apr. 2024. [Online].
Available: \url{https://www.ibm.com/es-es/think/topics/api}.
[Accessed: Jun. 5, 2026].

\bibitem{apirest}
R. T. Fielding, ``Architectural Styles and the Design of Network-based
Software Architectures,'' Ph.D. dissertation, Univ. of California,
Irvine, CA, USA, 2000. [Online].
Available: \url{https://ics.uci.edu/~fielding/pubs/dissertation/rest_arch_style.htm}.

\bibitem{websocket}
I. Fette and A. Melnikov, ``The WebSocket Protocol,''
RFC 6455, Internet Engineering Task Force (IETF), Dec. 2011.
[Online]. Available: \url{https://www.rfc-editor.org/rfc/rfc6455}.

\bibitem{tcp}
J. Postel, ``Transmission Control Protocol,''
RFC 793, Internet Engineering Task Force (IETF), Sep. 1981. [Online].
Available: \url{https://www.rfc-editor.org/rfc/rfc793}.

\bibitem{udp}
J. Postel, ``User Datagram Protocol,''
RFC 768, Internet Engineering Task Force (IETF), Aug. 1980. [Online].
Available: \url{https://www.rfc-editor.org/rfc/rfc768}.

\bibitem{http}
R. Fielding et al., ``HTTP Semantics,''
RFC 9110, Internet Engineering Task Force (IETF), Jun. 2022. [Online].
Available: \url{https://www.rfc-editor.org/rfc/rfc9110}.

\bibitem{tls}
E. Rescorla, ``The Transport Layer Security (TLS) Protocol Version 1.3,''
RFC 8446, Internet Engineering Task Force (IETF), Aug. 2018. [Online].
Available: \url{https://www.rfc-editor.org/rfc/rfc8446}.

\bibitem{bloqueo_http}
M. West, ``Mixed Content,''
W3C Candidate Recommendation, Aug. 2016. [Online].
Available: \url{https://www.w3.org/TR/mixed-content/}.

\bibitem{mqtt}
OASIS Standard, ``MQTT Version 5.0,''
OASIS, Mar. 7, 2019. [Online].
Available: \url{https://docs.oasis-open.org/mqtt/mqtt/v5.0/mqtt-v5.0.html}.

\bibitem{webrtc}
C. Jennings, H. Boström, and J.-M. Bruaroey, ``WebRTC 1.0:
Real-Time Communication Between Browsers,''
W3C Recommendation, Jan. 26, 2021. [Online].
Available: \url{https://www.w3.org/TR/webrtc/}.

\bibitem{signalling}
WebRTC.org, ``Peer Connections,''
Google. [Online].
Available: \url{https://webrtc.org/getting-started/peer-connections}.
[Accessed: Jun. 5, 2026].

\bibitem{json}
Ecma International, ``The JSON Data Interchange Format,''
ECMA-404, 2nd ed., Dec. 2017. [Online].
Available: \url{https://www.json.org/json-es.html}.







%% TFGs previos

\bibitem{pinon2024}
Piñón Mediero, Jorge,
\textit{Development of an attitude determination and control simulator for small satellites},
2024.
Treball de final de màster, Universitat Politècnica de Catalunya,
PRISMA-185586. URI: \url{http://hdl.handle.net/2117/410638}

\end{thebibliography}

\input{}
%%%%%%%%%%%%%%%%%%%%%%%%%%%%%%%%%%%%%%%%%%%%%%%%%%%%%%%%%%%%%%%%%%%%%%%%%%
%%%%%%                           APENDIXS                         %%%%%%%%
%%%%%%%%%%%%%%%%%%%%%%%%%%%%%%%%%%%%%%%%%%%%%%%%%%%%%%%%%%%%%%%%%%%%%%%%%%
\pagestyle{empty}  % no tocar

%% Descomentar una de les dues línies següents, en funció de:
%%  a) els apendixs s'encuadernaran apart (amb portada) 
%%  b) els apendixs s'enquadernen amb el mateix projecte (sense portada). 
%% Recordeu que si tot el document (amb apèndixs) excedeix les 100 pagines 
%% s'ha d'enquadernar a part
%\appendix\ambportada
\appendix\senseportada


%%%%%%%%%%%%%%%%%%%%%%%%%%%%%%%%%%%%%%%%%%%%%%%%%%%%%%%%%%%%%%%%%%%%%%%%%%
%%%%%% INCLOURE A PARTIR D'AQUI TOTS ELS CAPÍTOLS DELS APENDIXS   %%%%%%%%
%%%%%%%%%%%%%%%%%%%%%%%%%%%%%%%%%%%%%%%%%%%%%%%%%%%%%%%%%%%%%%%%%%%%%%%%%%

\input{appendix/AppendixA}
%\input{appendix/AppendixB}


%%%%%%%%%%%%%%%%%%%%%%%%%%%%%%%%%%%%%%%%%%%%%%%%%%%%%%%%%%%%%%%%%%%%%%%%%%
%%%%%%%%%%%%%%%%%%%%%%%%%%%%%%%%%%%%%%%%%%%%%%%%%%%%%%%%%%%%%%%%%%%%%%%%%%
%%%%%%%%%%%%%%%%%%%%%%%%%%%%%%%%%%%%%%%%%%%%%%%%%%%%%%%%%%%%%%%%%%%%%%%%%%
% i  aixo es tot! ;)
\end{document}






